\documentclass[11pt,a4paper]{article}
\usepackage[T1]{fontenc}
\usepackage[utf8]{inputenc}
\usepackage{amsmath,amsthm,mathtools}
\usepackage{newtxtext,newtxmath}
\usepackage[margin=27mm,headheight=15pt]{geometry}
\usepackage{microtype}
\usepackage{xcolor,booktabs,longtable,array,enumitem,fancyhdr}
\usepackage[most]{tcolorbox}
\usepackage{xurl}
\usepackage{hyperref}
\usepackage[nameinlink,noabbrev]{cleveref}
\definecolor{ink}{HTML}{17354A}
\definecolor{accent}{HTML}{176B75}
\definecolor{pale}{HTML}{F1F6F7}
\hypersetup{colorlinks=true,linkcolor=ink,citecolor=accent,urlcolor=accent,
 pdftitle={Self-Embeddings of the Surreal Numbers},
 pdfauthor={Research study prepared for Vladimir Reshetnikov}}
\setlength{\parskip}{4pt}
\setlength{\tabcolsep}{2.5pt}
\setlength{\emergencystretch}{2em}
\setlength{\parindent}{1.2em}
\setlist{itemsep=3pt,topsep=5pt}
\pagestyle{fancy}
\fancyhf{}
\fancyhead[L]{\small\color{ink}Self-Embeddings of the Surreal Numbers}
\fancyhead[R]{\small 3 October 2026}
\fancyfoot[C]{\thepage}
\renewcommand{\headrulewidth}{0.3pt}
\numberwithin{equation}{section}
\theoremstyle{plain}
\newtheorem{theorem}{Theorem}[section]
\newtheorem{proposition}[theorem]{Proposition}
\newtheorem{lemma}[theorem]{Lemma}
\newtheorem{corollary}[theorem]{Corollary}
\theoremstyle{definition}
\newtheorem{definition}[theorem]{Definition}
\newtheorem{example}[theorem]{Example}
\newtheorem{question}[theorem]{Research question}
\theoremstyle{remark}
\newtheorem{remark}[theorem]{Remark}
\newcommand{\No}{\mathbf{No}}
\newcommand{\Oz}{\mathbf{Oz}}
\newcommand{\On}{\mathbf{On}}
\newcommand{\R}{\mathbb{R}}
\newcommand{\Q}{\mathbb{Q}}
\newcommand{\Z}{\mathbb{Z}}
\newcommand{\N}{\mathbb{N}}
\newcommand{\id}{\operatorname{id}}
\newcommand{\supp}{\operatorname{supp}}
\newcommand{\Fix}{\operatorname{Fix}}
\newcommand{\im}{\operatorname{im}}
\newcommand{\bday}{\operatorname{b}}
\newcommand{\NF}[1]{\R((\omega^{#1}))}
\newcommand{\pref}{\preccurlyeq_s}
\newcommand{\spref}{\prec_s}
\newcommand{\concat}{\mathbin{\frown}}
\newcommand{\val}{\operatorname{v}}
\newcommand{\OrdExp}{\operatorname{Exp}}
\newcommand{\floorOz}[1]{\left\lfloor #1\right\rfloor_{\Oz}}
\newcommand{\code}[1]{\texttt{\small #1}}
\newtcolorbox{overview}{colback=pale,colframe=accent,boxrule=0.6pt,arc=1mm,
 left=3mm,right=3mm,top=2mm,bottom=2mm,breakable}
\newcolumntype{P}[1]{>{\raggedright\arraybackslash}p{#1}}
\begin{document}
\begin{titlepage}
\vspace*{13mm}
{\color{accent}\large RESEARCH ARTICLE}\par
\vspace{7mm}
{\color{ink}\Huge\bfseries Self-Embeddings of\\[3mm] the Surreal Numbers}\par
\vspace{7mm}
{\LARGE Rigidity, Hahn lifts, topology,\\and size-safe foundations}\par
\vspace{12mm}
{\large Prepared for Vladimir Reshetnikov}\par
{\large 3 October 2026}\par
\vspace{12mm}
\begin{overview}
\textbf{Central separation.} For every set $A\subseteq\No$, there is a
\emph{proper, cofinal, continuous ordered-field self-embedding} fixing
$A\cup\R\cup\On$ pointwise. It preserves all admissible Hahn sums,
normal-form truncations, omnific integrality, and omnific floors. Its image
is a closed nowhere-dense subfield. These constraints therefore do not
force surjectivity or rigidity.

In sharp contrast, preserving birthdays without increase forces an order
embedding to be the identity; an order embedding with simplicity-initial
image must be onto; and an exponential field self-embedding fixing the
natural valuation pointwise must be the identity.
\end{overview}
\vfill
\noindent\textbf{Repository basis.} \code{VladimirReshetnikov/ProveIt}, snapshot\par
\noindent\code{b0d4af4f30b24529fec877e2ad4a5952888146e9}.\par
\medskip
\noindent\small This is a research synthesis with proofs supplied in the text,
not a claim of publication priority. It distinguishes established literature,
repository results, deductions proved here, and unresolved directions.
No new Lean verification or independent repository build is claimed.
\end{titlepage}

\begin{abstract}
A self-embedding of the surreal numbers can mean an injection of sign
sequences, an order embedding, a simplicity-tree embedding, a field
embedding, an elementary embedding in a specified language, or a map
respecting valuation, normal forms, exponentiation, a derivation, or
particular distinguished subclasses. These notions have markedly different
rigidity and size behavior. We organize them into a common framework.
Our main constructive device is a two-stage Hahn lift that faithfully
converts any increasing class embedding of $\No$ into a strongly
$\R$-linear ordered-field self-embedding. It gives explicit proper maps
fixing every real, every ordinal, and any prescribed set of additional
surreals; it also gives bounded elementary subfield copies fixing any
prescribed set. We classify monomial strong embeddings by an additive
exponent embedding and a positive real character, compute their fixed
fields and images, and analyze their iterated-image cores.
For the full surreal order topology we prove a field-embedding dichotomy:
cofinal images give continuous embeddings, whereas noncofinal images are
discrete and the maps are nowhere continuous. We separate this topology
from Hahn summability, simplicity limits, and differentiation.
Building on the repository's displacement argument, we prove exponential
valuation rigidity without surjectivity and a uniqueness theorem for
cofinal exponential embeddings with the same value-group action.
Finally we give class-theoretic, universe-indexed, and Lean-oriented
formulations that avoid powerclass, truth-predicate, and improper-cut
paradoxes, and formulate a focused research program.
\end{abstract}

{\small\setlength{\parskip}{0pt}\tableofcontents}
\newpage

\section{The question is a family of questions}\label{sec:scope}

An embedding is always relative to a signature and to a category of
allowable maps. In this article, an \emph{order embedding} preserves and
reflects the numerical order; a \emph{field embedding} is unital; and a
\emph{proper embedding} is injective but not surjective. Unless a cutoff or
universe is explicitly named, the domain and codomain are the entire
proper class $\No$ in one fixed ambient set/class universe.

We use $\N=\{0,1,2,\ldots\}$.
The two fundamental orders must not be conflated. Numerical order is
linear. Simplicity is the tree order of being an initial segment of a sign
sequence. Being born earlier is weaker than being a proper sign prefix.
Likewise, an image that is closed under deleting late \emph{normal-form
terms} need not be closed under deleting late \emph{signs}.

\begin{longtable}{P{31mm}P{54mm}P{65mm}}
\toprule
Structure retained & Meaning & Main outcome here \\
\midrule
\endhead
Bare class carrier & Injective class map & Numerous proper maps; little arithmetic content. \\
Numerical order & $x<y\iff jx<jy$ & Copies in any open interval; proper cofinal copies too. \\
Simplicity tree & Prefix preservation/reflection & Rich tree maps; branch swaps are allowed without order. \\
Order and simplicity & Both relations & Proper substructures exist; automorphisms are rigid. \\
Birthdays or ambient cuts & Nonincreasing birthdays, or canonical cut preservation & Identity. \\
Additive group & Additive, ordered & Hahn relabeling of exponents is sufficient. \\
Ordered field & $0,1,+,\cdot,< $ & Proper cofinal and bounded copies exist. \\
Hahn structure & Strong sums and monomial data & Explicit classification within the monomial class. \\
Natural valuation & Comparisons, or pointwise values & These are different requirements; neither alone gives general field rigidity. \\
$\omega$-map / exponential & $j(\omega^x)=\omega^{jx}$ / $j(\exp x)=\exp(jx)$ & Distinct operations; valuation-pointwise compatibility gives rigidity. \\
Elementary structure & All first-order formulas of a fixed language & Automatic for ordered-field embeddings in the field language, not in arbitrary expansions. \\
Definability & Definable graph in a named language & O-minimal definable field self-embeddings are the identity. \\
Omnific / differential data & $\Oz$, its floor, a chosen derivation & Extra structures require separate compatibility statements. \\
\bottomrule
\end{longtable}

No finite list can exhaust all expansions of $\No$ by arbitrarily chosen
predicates. The meaningful aim is to cover the standard structures and
provide a method for analyzing additional constraints.

\subsection{Relation to existing work and to the repository}
The language of \emph{surreal substructures}, developed systematically by
Bagayoko and van der Hoeven, studies images isomorphic to
$(\No,<,\spref)$ and provides much more than the elementary prefix examples
below \cite{BVH}. Our constructions in the field category use the classical
normal-form and ordered-field theory of Conway and Gonshor, with the
set-sized field analysis of van den Dries and Ehrlich
\cite{Conway,Gonshor,vdDE}.

Kaplan, Krapp, and Serra's version~3 of \emph{Decomposing the automorphism
group of the surreal numbers} is the relevant recent comparison point
\cite{KKS}. It separates valuation, monomial, strong-linearity, exponential,
and $\omega$-map constraints. Its Question~5.4 asks about exponential
1-automorphisms; Question~5.6 concerns $\omega$-map-preserving field
automorphisms. The repository's exponential-rigidity report supplies a
finite displacement argument directed at the former question
\cite{RepoExp}. We reproduce the applicable proof and extend the
comparison argument to cofinal self-embeddings, while making no priority
claim and not treating the latter question as solved.

\begin{overview}
\textbf{Evidence boundary.} The pinned repository contains actual Lean
source for universe-indexed sign sequences and generic exponential-profile
lemmas. The inspected generic exponential module explicitly lists the
instantiation with the actual surreal exponential as pending. This article
uses conventional mathematical proofs for that bridge. The source
inspection is not an independent kernel build. The new Hahn-lift,
stabilizer, topology, and iteration arguments below are proved in this
article, not certified by the included finite regression tests.
\end{overview}

\section{A foundation that does not manufacture paradoxes}\label{sec:foundations}

\subsection{Set-length signs and set-supported normal forms}
Represent a surreal by a function $x:\alpha\to\{-,+\}$, where $\alpha$ is an
ordinal and hence a set. Its birthday is $\bday(x)=\alpha$. For comparison,
extend the sign function by the termination symbol $0$, ordered by
$-<0<+$. The first disagreement gives numerical order. Write
$x\pref y$ when $x$ is a sign prefix of $y$, and $x\spref y$ when it is a
proper prefix. The empty sequence is $0$.

The canonical option sets are
\begin{align}
 L_x&=\{x\restriction\beta:\beta<\bday(x),\ x(\beta)=+\},\\
 R_x&=\{x\restriction\beta:\beta<\bday(x),\ x(\beta)=-\}.
\end{align}
They are sets, and $x=\{L_x\mid R_x\}$. More precisely, the numbers between
these canonical options are exactly the sign extensions of $x$. For
arbitrary sets $L<R$, the notation $\{L\mid R\}$ denotes the unique simplest
surreal strictly between them. These standard sign and cut facts are the
foundational input used below \cite{Conway,Gonshor}.

Every surreal has a unique normal form
\begin{equation}\label{eq:nf}
 x=\sum_{a\in A}r_a\omega^a,
 \qquad r_a\in\R\setminus\{0\},
\end{equation}
where $A$ is a \emph{set} reverse well-ordered by numerical order. Equivalently,
its exponents form a strictly decreasing ordinal-indexed sequence. The
notation $\NF{E}$ will mean the class of such series with support contained
in $E$. It is only an additive subgroup when $E$ is an arbitrary subclass;
if $E$ is an ordered additive subgroup, it is a Hahn subfield.

For $x\ne0$, let $a_0=\max\supp(x)$ and set
\begin{equation}\label{eq:valuation}
 \val(x)=-a_0,\qquad \val(0)=\infty.
\end{equation}
Thus $\val(\omega^a)=-a$. The valuation group is identified with
$(\No,+,<)$, but the occurrence of $\No$ as value group should remain
logically distinct from its occurrence as field.

\begin{definition}[Strong summability]
A \emph{set-indexed} family $(x_i)_{i\in I}$ is Hahn-summable if the union of
its normal-form supports is reverse well-ordered and every exponent
occurs in only finitely many supports. The sum is then formed coefficient
by coefficient. A map is strongly additive if it preserves these families
and their sums; it is strongly $\R$-linear if it also respects real scalar
multiplication.
\end{definition}
This is a support condition, not a limit in the full surreal order topology.
In particular, an arbitrary proper-class family is not a legal Hahn sum.

\subsection{Three legitimate size registers}
\paragraph{Definable classes in ZFC.}
Each surreal is a set code, and a specified embedding is given by a formula,
possibly with set parameters. Assertions about every definable embedding
are metatheoretic schemes, not quantification over a set of all formulas
with a universally available truth predicate. All explicit maps in our
Hahn constructions are definable from their stated data.

\paragraph{Two-sorted class theory.}
In GBC, or in a suitable NBG presentation, an embedding may be a class graph.
Class replacement ensures that its image of a set is a set. This suffices
for mapping supports, taking preimages of a set inside an injective image,
and performing the explicit constructions below. Global choice is useful
for a global back-and-forth using a set-like enumeration; it is not needed
for the displayed rational formulas and Hahn relabelings.

One may quantify over an individual class embedding, or finitely many
class embeddings, in class theory. One may \emph{not} form a class whose
members are all those proper-class graphs: classes are not elements of
classes. Thus ``the group of all automorphisms of $\No$'' is an external
collection with operations, or an abbreviation for finite class-quantified
statements. A genuine carrier of all such maps requires a higher sort,
a larger universe, or a restricted set-coded family. KKS explicitly
highlights this point in its foundational conventions \cite{KKS}.

\paragraph{Universe-indexed types.}
Fix a small universe $U$ and work externally in a larger universe $V$.
The surreals with $U$-small ordinal lengths form a $V$-set, denoted
$\No_U$. They are not a $U$-set. Maps $\No_U\to\No_U$, maps
$\No_U\to\No_V$, and restrictions of maps on a still larger surreal
universe are different objects. An external set of signs need not be an
internally available support. An external cut need not be an admissible
internal cut.

This is directly reflected by the inspected repository declaration
\begin{verbatim}
structure SignSequence : Type (u + 1) where
  birthday : Ordinal.{u}
  signAt : Ordinal.{u} -> SignType
  signAt_ne_zero_iff : forall i, signAt i != 0 <-> i < birthday
\end{verbatim}
The display uses ASCII logical notation for readability; the exact Lean
source uses its usual Unicode syntax. The file proves that bounded birthday
fragments are small and that the entire carrier is not small at level $u$
\cite{RepoSigns}.

\subsection{What is forbidden, and what needs extra axioms}
The expression $\{\No\mid\varnothing\}$ is not a surreal cut: its left side
is a proper class. Nor is there a single surreal sign sequence of length
``all ordinals.'' These errors would falsely produce an element above every
surreal and immediately contradict ordinary field arithmetic.

Similarly, a universal satisfaction predicate for an arbitrary proper-class
structure is not supplied merely by writing the word ``elementary.'' One
can state preservation for each fixed formula. In the ordered-field
language, quantifier elimination gives an especially clean reduction to
finite algebraic/order data.

Set-valued recursion along $\On$, with a set of earlier values at each set
stage, is different from recursion that assigns an entire class at each
stage. The latter leads to elementary transfinite recursion principles
stronger than bare GBC in the usual class-theoretic setting
\cite{HamkinsETR}. We do not silently use such principles to construct the
maps in this article. In particular, a set-like global back-and-forth has
set-sized partial maps at every ordinal stage; its final union is a class,
not a set.

Finally, a field self-embedding of $\No$ is not an elementary embedding of
the set-theoretic universe $(V,\in)$. The underlying membership relation on
set codes is not part of the ordered-field signature. Large-cardinal and
Kunen-style questions about universe embeddings therefore cannot be
transferred to the present field maps without a separate interpretation
theorem. Coding $\No$ by sets does not provide such a theorem by itself.

\section{Order and simplicity: flexibility and rigidity}\label{sec:tree}

\subsection{Three explicit order embeddings}
\begin{proposition}[Interval compression]\label{prop:compression}
The map
\[
 c(x)=\frac{x}{1+|x|}
\]
is an increasing bijection $\No\to(-1,1)$, with inverse
$c^{-1}(y)=y/(1-|y|)$. Consequently every nonempty open interval contains
an order-isomorphic copy of the entire surreal class.
\end{proposition}
\begin{proof}
For $x\ge0$, the formula is $1-(1+x)^{-1}$, which is strictly increasing
with image $[0,1)$. The negative half is obtained by oddness. The inverse
formula is verified separately on the two signs. An affine increasing
bijection carries $(-1,1)$ to any prescribed bounded open interval; the
half-line variants follow by equally elementary fractional formulas.
\end{proof}
No cardinal paradox occurs: an interval of a proper-class field is itself
a proper class. This map is not additive and is not a field embedding.

\begin{proposition}[Prefix copies]\label{prop:prefix}
For a fixed sign sequence $p$, put $P_p(x)=p\concat x$. Then $P_p$ preserves
and reflects order and simplicity. Its image is the convex class of all
extensions of $p$, and
\[
 \bday(P_p(x))=\bday(p)+\bday(x),
\]
where addition is ordinal addition. If $p\ne0$, the embedding is proper.
\end{proposition}
\begin{proof}
The common prefix $p$ postpones every first disagreement by its length.
It also preserves exactly the initial-segment relations among the suffixes.
The characterization of the image is immediate. If two extensions bracket
another sign sequence numerically, that sequence cannot disagree with $p$
at an earlier position, so the image is convex. The empty sequence is not
in the image when $p$ is nonempty.
\end{proof}
Ordinal addition matters: a nonempty prefix need not strictly increase an
infinite birthday. For example $1+\omega=\omega$.

\begin{proposition}[Sign repetition]\label{prop:repeat}
For an integer $m\ge2$, let $D_m(x)$ replace every sign of $x$ by $m$
consecutive copies of that sign. Then $D_m$ is a proper, cofinal and
coinitial order-and-simplicity embedding fixing $0$, with
\[
 \bday(D_m(x))=m\cdot\bday(x).
\]
Its image is discrete in the ambient order topology and $D_m$ is nowhere
continuous.
\end{proposition}
\begin{proof}
At a first differing sign, repeated minus, termination, and repeated plus
occur in the same order as the original three symbols. Prefix relations
are preserved and reflected. The one-sign sequence $+$ is omitted, so the
map is proper. Images of sufficiently long all-plus sequences lie above
any given surreal; all-minus sequences give coinitiality.

For $z=D_m(x)$, let $z^-=z\concat(-)$ and $z^+=z\concat(+)$. The interval
$(z^-,z^+)$ contains no other point of the image: any image extending $z$
has at least two copies of its next sign and lies beyond the corresponding
endpoint. All other image points disagree earlier or terminate earlier.
Thus each image point is isolated. The source has no isolated points, so
continuity at any point is impossible.
\end{proof}
For example $D_2(1)=2$, but $D_2(1/2)$ has signs $++--$ and equals $5/4$;
therefore $D_2$ is not additive. Also $D_m(\omega)=\omega$, since
$m\cdot\omega=\omega$. These examples already show that cofinality,
continuity, and birthday growth are different issues for order embeddings.

\subsection{Simplicity without order}
The rooted binary sign tree has many automorphisms if the order of the two
children is forgotten. At each node one may decide whether to interchange
the two outgoing branches, recursively. These automorphisms preserve
birthdays and prefix relations but need not preserve numerical order.
Global sign reversal $x\mapsto-x$ is the simplest example: it reverses
numerical order and is not a unital field homomorphism. Every order-reversing
self-embedding is the composite of negation with an order-preserving one.

\begin{lemma}[Reflection follows from order]\label{lem:reflection}
An order embedding that preserves sign-prefix relations automatically
reflects them.
\end{lemma}
\begin{proof}
If $y\spref x$, preservation and injectivity preclude $jx\pref jy$.
If $x,y$ are incomparable, their longest common prefix $p$ lies strictly
between them numerically, after interchanging $x,y$ if needed. The images
extend $jp$ and lie on opposite sides of it. They therefore depart along
opposite branches and are incomparable. These exhaust the alternatives
to $x\pref y$.
\end{proof}

\subsection{Rigidity statements with different hypotheses}
\begin{theorem}[Birthday rigidity]\label{thm:birthday}
If $j:\No\to\No$ is an order embedding and
$\bday(jx)\le\bday(x)$ for every $x$, then $j=\id$.
\end{theorem}
\begin{proof}
Induct on $\alpha=\bday(x)$. Every canonical option of $x$ has smaller
birthday and is fixed by induction. Therefore $jx$ lies between $L_x$ and
$R_x$, so it extends $x$. Its birthday cannot exceed $\alpha$ by hypothesis;
hence it is exactly $x$.
\end{proof}

\begin{proposition}\label{prop:treebirthday}
Every simplicity embedding satisfies $\bday(jx)\ge\bday(x)$.
Every order-and-simplicity automorphism of $\No$ is the identity.
\end{proposition}
\begin{proof}
The proper prefixes of $x$ form a chain of order type $\bday(x)$; their
images form a strictly increasing chain of proper prefixes of $jx$.
Consequently $\bday(x)\le\bday(jx)$. For an automorphism, apply the same
argument to its inverse to obtain equality of birthdays, then use
\cref{thm:birthday}.
\end{proof}
This is compatible with the proper maps $P_p$ and $D_m$: an embedding into
an image is not an automorphism of the whole tree.

\begin{proposition}[Canonical construction rigidity]\label{prop:cutrigid}
An order embedding preserving ambient canonical cuts,
\[
 j\{L\mid R\}=\{j[L]\mid j[R]\}
\]
for every admissible set cut, is the identity. It is enough to require this
for the canonical cuts of individual surreals.
\end{proposition}
\begin{proof}
Transfinite induction fixes all options and then fixes the simplest
number between them. The empty cut fixes $0$ at the first step.
\end{proof}
Equivalently, preserving the root, the two labeled immediate successors,
and unions of all set-sized prefix chains forces the identity. Omitting
limit-chain preservation leaves room for extra signs at limit birthdays.
Preservation only for countable chains does not justify a conclusion at
arbitrary uncountable birthdays.

\begin{remark}[Relative cuts are not ambient cuts]
If $S$ is the image of an order-and-simplicity embedding, the image of a cut
is the simplest suitable point \emph{inside $S$}. It need not be the
simplest suitable point in all of $\No$. This relative notion is central
to surreal substructures \cite{BVH}; replacing it with the ambient notion
would incorrectly rule out the elementary prefix examples.
\end{remark}


\subsection{Closed substructures and their fixed points}
The rigidity of automorphisms implies uniqueness of the parameterization
of any given surreal substructure: compose two proposed parameterizations
with one another's inverses. Monomials and purely infinite numbers supply
further standard substructure examples beyond the prefix cones; their
substructure parameterizations are not thereby field embeddings
\cite{Gonshor,BVH}.

There is also a useful fixed-point theorem under an explicit closure
hypothesis. The following gives a direct proof of this basic closed-
substructure phenomenon, part of the broader theory in \cite{BVH}.

\begin{theorem}[Closed sign images have a full fixed-point substructure]
\label{thm:closedtree}
Let $j$ be an order-and-simplicity embedding. Suppose its image is closed
under unions of nonempty set-sized sign-prefix chains. Then $j$ commutes
with those unions, and $\Fix(j)$ is itself a surreal substructure.
\end{theorem}
\begin{proof}
Let $C$ be a nonempty set-sized prefix chain and $x=\bigcup C$.
By the image-closure assumption, $y=\bigcup j[C]$ lies in the image.
Its preimage is an upper bound of $C$, while $jx$ is an upper bound of
$j[C]$. The two least-upper-bound properties give $jx=y$.

Whenever $u\pref ju$, put
\[
 c_j(u)=\bigcup_{n<\omega}j^n(u).
\]
This is fixed by union preservation. It is the least fixed sign extension
of $u$: if $u\pref z=jz$, induction gives $j^n u\pref z$ for every $n$.
In particular $c_j(0)$ is the least fixed point in simplicity.
If $x$ is fixed, then $u=x\concat(\pm)$ satisfies $u\pref ju$.
Indeed $ju$ extends $x$ and lies on the same numerical side of $x$ as $u$.
Thus $c_j(x\concat(-))$ and $c_j(x\concat(+))$ are the least fixed
extensions on the two respective branches.

Construct a map $\Phi$ from sign sequences by sending the empty sequence
to $c_j(0)$, adjoining a sign by the corresponding least-fixed-extension
operation, and taking unions at limit lengths. All these unions are sets;
union preservation by $j$ keeps every output fixed. At each successor
step the image strictly extends its predecessor on the specified side.
Hence $\Phi$ preserves and reflects numerical order and simplicity.

To see that its image is all of $\Fix(j)$, take a fixed $z$. Its least
fixed prefix is $c_j(0)$. As long as the constructed prefix is not $z$,
follow the next sign of $z$ and then take the least fixed extension on
that branch. This remains a prefix of $z$. At limits take unions. The
process must reach $z$: otherwise it would produce more than
$\bday(z)$ strict prefix extensions inside the set-length sequence $z$.
Thus $\Phi$ is onto the fixed-point class.
\end{proof}

For example, $D_m$ commutes with sign-chain unions and its image is closed
under them, so its fixed points form a surreal substructure. Yet $D_m$
is nowhere continuous in the numerical order topology by
\cref{prop:repeat}. The two kinds of continuity are not interchangeable.

\section{Initial images and the distinction between two truncations}\label{sec:initial}

\begin{definition}
A subclass $S\subseteq\No$ is \emph{simplicity-initial} if $x\in S$ and
$y\pref x$ imply $y\in S$. A normal-form truncation discards a final segment
of the decreasing list of exponents in \eqref{eq:nf}. These two operations
are unrelated definitions and must be tracked separately.
\end{definition}

\begin{theorem}[An initial full-sized order copy is the whole class]
\label{thm:initial}
If an order embedding $j:\No\to\No$ has simplicity-initial image $S$, then
$S=\No$. Simplicity preservation by $j$ is not required.
\end{theorem}
\begin{proof}
The order $S$ inherits interpolation for every set cut: take preimages
under $j$, interpolate in $\No$, then apply $j$.
Suppose $S\ne\No$, and choose an omitted $x$ of least birthday. Every
proper prefix of $x$ belongs to $S$. In particular $L_x,R_x\subseteq S$.
Interpolation in $S$ gives $y\in S$ with $L_x<y<R_x$. Such a $y$ extends
$x$. Initiality then gives $x\in S$, a contradiction.
\end{proof}
The conclusion is surjectivity, not identity. A nonidentity ordered-field
automorphism has image $\No$, which is certainly initial. Adding simplicity
preservation gives identity by \cref{prop:treebirthday}.

This theorem does not conflict with initial embeddings of set-sized
ordered fields into $\No$ \cite{EhrlichKaplan}. A set-sized field need not
interpolate every set cut of its underlying order. The full class $\No$
does, and that extra property drives the proof.

\begin{proposition}[A proper subfield is not convex]\label{prop:convexfield}
A convex subfield of an ordered field is the whole field. More generally,
a subfield containing a nonempty open interval is the whole field.
\end{proposition}
\begin{proof}
For the first statement, if $a>0$ lies above the convex subfield, then
$0<a^{-1}<1$ puts $a^{-1}$ in the subfield, hence also $a$.
For the second, subtract two points of the interval to obtain an interval
$(-\delta,\delta)$ inside the subfield. For arbitrary $x$ in the ambient
field choose $0<t<\delta/(1+|x|)$. Both $t$ and $tx$ lie in that interval,
so $x=(tx)/t$ belongs to the subfield.
\end{proof}
Thus the convex prefix copies and interval copies are inherently different
from proper field copies. The field images constructed next are closed
under normal-form truncation, but by \cref{thm:initial} they cannot be
simplicity-initial.

\section{Hahn relabeling and a double-lift construction}\label{sec:hahn}

The principal construction is deliberately modular. Relabeling the
exponents of a surreal needs only an order embedding and produces an
\emph{additive} embedding. Relabeling exponents by an \emph{additive}
embedding then produces a field embedding. Keeping these two levels
separate prevents a common false claim that every monotone change of
exponent preserves multiplication.

\subsection{The additive lift of an order embedding}
Let $f:\No\to\No$ be an increasing class embedding. Define
\begin{equation}\label{eq:addlift}
 A_f\left(\sum_a r_a\omega^a\right)
       =\sum_a r_a\omega^{f(a)}.
\end{equation}
The sum on the right is legal: the image of its support is a set and an
order-isomorphic copy of a reverse well-order.

\begin{theorem}[Additive lift]\label{thm:addlift}
The map $A_f$ is a strongly $\R$-linear ordered additive-group embedding.
It preserves and reflects Hahn summability, and
\begin{align}
 \im A_f&=\NF{f(\No)},\label{eq:addimage}\\
 A_f\circ A_g&=A_{f\circ g},\label{eq:addfunctor}\\
 \Fix(A_f)&=\NF{\Fix(f)}.\label{eq:addfixed}
\end{align}
It is surjective exactly when $f$ is surjective, and its image is cofinal
exactly when $f(\No)$ is cofinal.
\end{theorem}
\begin{proof}
Addition and real scalar multiplication are coefficientwise. An increasing
injection preserves the largest differing exponent and its coefficient,
so it preserves order and is injective. The two summability conditions are
invariant under an injective order relabeling: support order type is
unchanged, as is the number of indices contributing to a coefficient.
Equation~\eqref{eq:addimage} follows by pulling back a support contained in
$f(\No)$ under the inverse class map on that image. Class replacement
makes the pulled-back support a set. Composition is immediate termwise.

If $A_f(x)=x$, its support $B$ satisfies $f[B]=B$. The restriction of $f$ is
an order automorphism of the reverse well-order $B$, hence is the identity
on $B$. Conversely, support in $\Fix(f)$ gives a fixed series. This proves
\eqref{eq:addfixed} without assuming $f$ is onto on the whole class.

Surjectivity follows from the image formula; an omitted exponent omits its
monomial. If $f(\No)$ is cofinal, images of monomials dominate any given
normal form. If it is not cofinal, choose a strict upper bound $b$ for it.
Every nonzero $A_f(x)$ has largest exponent below $b$, so
$|A_f(x)|<\omega^b$. Its image is bounded.
\end{proof}

Even the image of $1$ changes in general: $A_f(1)=\omega^{f(0)}$.
The symbol $A_f$ should never be called a field embedding unless the
additional multiplicative conditions have been verified.

\subsection{The field lift of an additive embedding}
Let $\tau:(\No,+,<)\to(\No,+,<)$ be an ordered additive-group embedding.
Define
\begin{equation}\label{eq:fieldlift}
 H_\tau\left(\sum_a r_a\omega^a\right)
       =\sum_a r_a\omega^{\tau(a)}.
\end{equation}

\begin{theorem}[Field lift]\label{thm:fieldlift}
The map $H_\tau$ is a unital ordered-field embedding, fixes $\R$ pointwise,
and is strongly $\R$-linear. It preserves and reflects Hahn summability.
Moreover,
\begin{align}
 \im H_\tau&=\NF{\tau(\No)},\label{eq:fieldimage}\\
 H_\tau H_\rho&=H_{\tau\circ\rho},\label{eq:fieldfunctor}\\
 \val(H_\tau x)&=\tau(\val x)\quad(x\ne0),\label{eq:fieldvalue}\\
 \Fix(H_\tau)&=\NF{\Fix(\tau)}.\label{eq:fieldfixed}
\end{align}
Surjectivity and cofinality are respectively equivalent to surjectivity
and cofinality of $\tau$.
\end{theorem}
\begin{proof}
The additive-lift proof applies to addition and supports. Since
$\tau(0)=0$, constants are fixed. For multiplication, the identity
$\tau(a+b)=\tau(a)+\tau(b)$ carries every convolution term
$r_as_b\omega^{a+b}$ to the corresponding term of the product of the
images. The usual finiteness of contributions to a Hahn product coefficient
is unchanged. Thus multiplication is preserved. Injectivity then implies
preservation of inverses, because $H_\tau(x)H_\tau(x^{-1})=1$.

The image, composition, fixed-point, and cofinality statements follow by
the same support arguments as above. Finally, if $a$ is the leading
exponent of $x$, then
$\val(H_\tau x)=-\tau(a)=\tau(-a)=\tau(\val x)$.
\end{proof}

The image is a real-closed field, because it is isomorphic to $\No$. This
uses the real closedness of the surreal field, not an unsupported assertion
that arbitrary birthday cutoffs are fields. Each image is also closed
under all normal-form truncations.

\begin{theorem}[Double lift]\label{thm:doublelift}
For an increasing class embedding $f:\No\to\No$, put
\begin{equation}\label{eq:doublelift}
 J_f=H_{A_f}.
\end{equation}
Then $J_f$ is a strongly $\R$-linear ordered-field self-embedding, and
\[
 J_fJ_g=J_{f\circ g},\qquad J_{\id}=\id.
\]
The assignment is faithful: $J_f=J_g$ implies $f=g$. It preserves and
reflects surjectivity and cofinality.
\end{theorem}
\begin{proof}
The first lift is an ordered additive embedding, so the second is a field
embedding. Both composition laws give the displayed identity. Faithfulness
can be read off from a single family of test elements:
\begin{equation}\label{eq:testelement}
 J_f\bigl(\omega^{\omega^a}\bigr)=\omega^{\omega^{f(a)}}.
\end{equation}
Injectivity of the $\omega$-map twice recovers $f(a)$. The remaining
assertions follow by applying the image criteria for the two lifts.
\end{proof}

The construction converts an order-embedding problem into a field-embedding
problem without pretending that the original $f$ was additive. Statements
about the ``monoid'' of all such maps are understood in the external or
finite-class-quantified sense of \cref{sec:foundations}. The formulas also
give ordinary monoid homomorphisms on any legitimate set-coded family
closed under composition.

\subsection{Classification with real weights}
There is a useful exact classification within a specified strong monomial
class. It should not be mistaken for a classification of every field
embedding of $\No$.

\begin{theorem}[Weighted monomial classification]\label{thm:weighted}
Let $\tau$ be an ordered additive embedding of $\No$, and let
$\chi:(\No,+)\to(\R_{>0},\cdot)$ be a group homomorphism. Then
\begin{equation}\label{eq:weighted}
 H_{\tau,\chi}\left(\sum_a r_a\omega^a\right)
 =\sum_a r_a\chi(a)\omega^{\tau(a)}
\end{equation}
is a strongly $\R$-linear ordered-field embedding. Conversely, every
strongly $\R$-linear field embedding sending every monomial $\omega^a$ to
a positive real multiple of a monomial has this unique form.
Its image is $\NF{\tau(\No)}$, and
\begin{equation}\label{eq:weightedfixed}
 \Fix(H_{\tau,\chi})=
 \NF{\Fix(\tau)\cap\ker\chi}.
\end{equation}
Here $\ker\chi=\{a:\chi(a)=1\}$. Composition is given by
\[
 (\tau_1,\chi_1)(\tau_2,\chi_2)
 =\left(\tau_1\tau_2,
     a\longmapsto\chi_2(a)\chi_1(\tau_2(a))\right).
\]
\end{theorem}
\begin{proof}
The character law and additivity of $\tau$ prove multiplicativity on
monomials; strong sums extend it to all products. Positive coefficients
preserve the sign of the leading term. Conversely, unique normal forms of
the images of monomials determine $\tau(a)$ and $\chi(a)$. Multiplication
of monomials forces additivity and the character law. The strict order of
Archimedean classes of monomials forces $\tau$ to be an order embedding.
Strong real linearity determines the map on every series.

Every series supported on $\tau(\No)$ is in the image: pull back each
exponent and divide its real coefficient by the corresponding nonzero
$\chi(a)$. For fixed points, equality of supports again makes $\tau$ the
identity on the support. Equality of nonzero coefficients then requires
$\chi(a)=1$ at each exponent. The composition formula follows by applying
the two maps to $\omega^a$.
\end{proof}

General strong embeddings can send a monomial to a monomial times a
nonconstant unit, and general field embeddings need not be strong at all.
The theorem makes no claim about those larger classes. The full
surreal-field and Hahn-field automorphism literature includes such distinctions
\cite{KKS,KS}.

\section{Proper embeddings fixing every real and ordinal}\label{sec:stabilizers}

\subsection{A completely explicit example}
Define the increasing map
\begin{equation}\label{eq:fminus}
 f_-(a)=
 \begin{cases}
 \displaystyle\frac{a}{1-a},&a<0,\\[2mm]
 a,&a\ge0.
 \end{cases}
\end{equation}
It maps the negative half-line bijectively onto $(-1,0)$ and fixes the
nonnegative half-line. Therefore it is an order isomorphism
$\No\to(-1,\infty)$, proper and cofinal. Let
\[
 \tau_-=A_{f_-},\qquad J_-=H_{\tau_-}.
\]

\begin{theorem}[An ordinal-fixing proper field copy]\label{thm:explicit}
The field embedding $J_-$ is proper and cofinal, fixes every real and every
ordinal, and preserves all Hahn-summable families and all normal-form
truncations. It moves the explicit omnific integer
\[
 J_-\bigl(\omega^{\omega^{-1}}\bigr)
     =\omega^{\omega^{-1/2}}.
\]
The element $\omega^{\omega^{-1}}$ does not belong to its image.
\end{theorem}
\begin{proof}
Properness and cofinality follow from \cref{thm:doublelift}. Real constants
are fixed by every field lift. The Cantor normal form of an ordinal has
ordinal exponents and is its Conway normal form. Since $f_-$ fixes every
nonnegative exponent, $A_{f_-}$ fixes every ordinal. Applying the outer
lift now fixes the Cantor normal form of every ordinal as well.

The displayed motion follows from $f_-(-1)=-1/2$. The exponent $-1$ is
omitted by $f_-$, so $\omega^{-1}$ is omitted by $A_{f_-}$. Consequently the
outer monomial $\omega^{\omega^{-1}}$ is omitted by $J_-$.
\end{proof}

There is no contradiction between ``omitted'' and ``moved'': an element
of the domain need not belong to the range, although it certainly has an
image. Also, $J_-(\omega^{-1})=\omega^{-1}$, because $A_{f_-}$ fixes the
constant $-1$. The map does not merely replace every visible occurrence
of $\omega$ by a new base; the level of each normal-form exponent matters.

\subsection{Fixing an arbitrary prescribed set}
\begin{definition}[Two-level exponent footprint]
For a set $A\subseteq\No$, define the set
\begin{equation}\label{eq:footprint}
 E(A)=\bigcup_{x\in A}\ \bigcup_{a\in\supp(x)}\supp(a).
\end{equation}
This consists of the exponents occurring in the normal forms of the
exponents occurring in members of $A$.
\end{definition}
The expression is a set by replacement and union. That elementary size
observation makes the following stabilizer theorem possible.

\begin{lemma}[Footprint criterion]\label{lem:footprint}
If $f$ fixes $E(A)$ pointwise, then $J_f$ fixes $A$ pointwise.
\end{lemma}
\begin{proof}
For every outer exponent $a$ in a member of $A$, the support of $a$ is
fixed by $f$. Thus $A_f(a)=a$. The outer field lift then leaves every
normal form in $A$ unchanged.
\end{proof}

\begin{theorem}[Large stabilizers, including proper maps]\label{thm:stabilizer}
For every set $A\subseteq\No$, there exists a proper cofinal strongly
$\R$-linear field embedding fixing $A\cup\R\cup\On$ pointwise.
There also exists a nonidentity strongly $\R$-linear field automorphism
fixing that same class pointwise.
\end{theorem}
\begin{proof}
Choose $\beta<0$ below every member of the set $E(A)$, using set-cut
interpolation. Define
\begin{equation}\label{eq:fbeta}
 f_\beta(a)=
 \begin{cases}
 \displaystyle\beta+\frac{a-\beta}{1-(a-\beta)},&a<\beta,\\[2mm]
 a,&a\ge\beta.
 \end{cases}
\end{equation}
This is an increasing bijection from $\No$ onto $(\beta-1,\infty)$.
It fixes $E(A)$ and every nonnegative exponent. The footprint criterion
fixes $A$; the ordinal argument of \cref{thm:explicit} fixes $\On$; and
\cref{thm:doublelift} supplies properness and cofinality of $J_{f_\beta}$.

For the automorphism, replace the lower branch by an affine dilation:
\begin{equation}\label{eq:gbeta}
 g_\beta(a)=
 \begin{cases}
 \beta+2(a-\beta),&a<\beta,\\
 a,&a\ge\beta.
 \end{cases}
\end{equation}
This is a nonidentity order automorphism of $\No$ with the same fixed
upper tail. Its double lift is a nonidentity field automorphism by
faithfulness, with all the required fixed points.
\end{proof}

\begin{corollary}[No set of extra parameters is a rigidity base]
Even after all reals and all ordinals are fixed, no set of additional
surreal parameters forces every strongly $\R$-linear monomial field
self-embedding to be the identity, or even to be surjective.
\end{corollary}
The fixed class $\On$ is itself proper. The statement is therefore
stronger than simply observing that a saturated object has nontrivial
automorphisms over a small set.

\subsection{Bounded field copies fixing prescribed data}
\begin{theorem}[Bounded copies]\label{thm:boundedcopy}
For every set $A\subseteq\No$, there is a strongly $\R$-linear field
self-embedding fixing $A\cup\R$ whose entire image is bounded in $\No$.
It can also be chosen to fix $\omega$.
\end{theorem}
\begin{proof}
Choose $\gamma>0$ above every member of $E(A)$, and put
\begin{equation}\label{eq:fupper}
 f^+_\gamma(a)=
 \begin{cases}
 a,&a\le\gamma,\\[1mm]
 \displaystyle\gamma+\frac{a-\gamma}{1+(a-\gamma)},&a>\gamma.
 \end{cases}
\end{equation}
The range is $(-\infty,\gamma+1)$. Therefore for every $x$,
\[
 |A_{f^+_\gamma}(x)|<\omega^{\gamma+1}.
\]
The exponents of every image under $J_{f^+_\gamma}$ are consequently below
$\omega^{\gamma+1}$, giving the common strict bound
\begin{equation}\label{eq:boundedbound}
 |J_{f^+_\gamma}(y)|<\omega^{\omega^{\gamma+1}}
 \qquad(y\in\No).
\end{equation}
The footprint criterion fixes $A$. Constants are fixed, and
$f^+_\gamma(0)=0$ gives $A_{f^+_\gamma}(1)=1$, hence $J_{f^+_\gamma}(\omega)=\omega$.
\end{proof}
A bounded subfield is not an Archimedean subfield. It can contain infinite
elements, all their inverses, and arbitrarily complicated prescribed set
families, while a still larger ambient surreal bounds the entire class
image. Such an embedding cannot fix every ordinal, since the ordinals are
cofinal in $\No$.

\section{Omnific floors, fixed fields, and iteration}\label{sec:fixed}

\subsection{Omnific integrality is preserved and reflected}
Conway's omnific integers have the normal-form description
\begin{equation}\label{eq:oz}
 \Oz=\left\{x:\supp(x)\subseteq\No_{\ge0},\quad [x]_0\in\Z\right\},
\end{equation}
where $[x]_0$ is the coefficient of $\omega^0$, taken as zero when absent.
Every surreal has a unique omnific floor $a$ with $a\le x<a+1$.
These are the usual omnific conventions, also used in the repository's
constructed omnific ring \cite{Conway,RepoMain}.

\begin{proposition}[Omnific compatibility]\label{prop:oz}
Every weighted monomial embedding $H_{\tau,\chi}$ satisfies
\[
 x\in\Oz\iff H_{\tau,\chi}(x)\in\Oz,
 \qquad
 H_{\tau,\chi}(\floorOz{x})=
 \floorOz{H_{\tau,\chi}(x)}.
\]
It preserves the constant coefficient and commutes with normal-form
truncations at corresponding support positions.
\end{proposition}
\begin{proof}
An additive ordered embedding fixes zero and preserves and reflects the
sign of every exponent. Since $\chi(0)=1$, the constant coefficient is
unchanged. This proves the membership equivalence. If $a$ is the omnific
floor of $x$, apply the ordered field embedding to $a\le x<a+1$.
The image of $a$ is omnific, and uniqueness of the floor gives the formula.
Termwise relabeling plainly commutes with deleting a final support segment.
\end{proof}
This includes the negative-infinitesimal boundary: the floor of a negative
infinitesimal is $-1$, not zero. The proof does not rely on ignoring that
boundary. Preserving $\Oz$ as a predicate is much weaker than fixing it
pointwise: the moved element in \cref{thm:explicit} belongs to $\Oz$.

\subsection{Fixed fields do not determine surjectivity}
\begin{proposition}[General fixed-field facts]\label{prop:fixedgeneral}
The fixed field of any ordered-field self-embedding of $\No$ is real
closed. No point has a nontrivial finite orbit, and
$\Fix(j^n)=\Fix(j)$ for every positive integer $n$.
\end{proposition}
\begin{proof}
The fixed points form a subfield. For a polynomial with fixed coefficients,
its finite set of roots in $\No$ is carried injectively into itself. The
induced map is a strictly increasing permutation of a finite ordered set,
so fixes every root. In particular positive square roots and a root of
any odd-degree polynomial lie in the fixed field, proving real closedness.
If $jx>x$, its forward iterates are strictly increasing; if $jx<x$, they
are strictly decreasing. Neither case gives a periodic point.
\end{proof}

Let
\begin{equation}\label{eq:B}
 B=\NF{\No_{\ge0}}.
\end{equation}
This is the additive subgroup of purely infinite parts plus real constants.
It need not be confused with $\Oz$, whose constant term must be integral.
Since $\Fix(f_-)=\No_{\ge0}$, the two lift formulas give
\begin{equation}\label{eq:Jfixed}
 \Fix(\tau_-)=B,\qquad \Fix(J_-)=\NF{B}.
\end{equation}
Now take $g(a)=2a$ for $a<0$ and $g(a)=a$ for $a\ge0$. Its double lift is
an automorphism, but it has exactly the same fixed field $\NF{B}$ as the
proper embedding $J_-$. Consequently neither this fixed field nor its
inclusion of all ordinals can detect whether an embedding is onto.

\subsection{The largest reversible core}
\begin{definition}
For an injective field endomorphism $j$, define its iterated-image core by
\begin{equation}\label{eq:core}
 K_\infty(j)=\bigcap_{n<\omega}j^n(\No).
\end{equation}
This is an intersection specified by quantification over natural numbers,
not a set containing proper-class fields as elements.
\end{definition}

\begin{proposition}\label{prop:core}
The restriction of $j$ to $K_\infty(j)$ is an automorphism. This is the
largest subfield on which $j$ restricts to an automorphism. The core is
real closed and can strictly contain $\Fix(j)$.
\end{proposition}
\begin{proof}
Forward invariance is immediate. For $x\in K_\infty(j)$ let $j(y)=x$.
Since $x\in j^{n+1}(\No)$ for every $n$, injectivity implies
$y\in j^n(\No)$ for every $n$. Thus $y$ lies in the core, proving
surjectivity of the restriction. A subfield on which $j$ is bijective
lies in every iterated image, proving maximality.

Each iterated image is real closed. Their intersection is real closed:
for a polynomial over the intersection, all of its real roots in $\No$
belong to each image field, since a real-closed subfield has no proper
ordered algebraic extension inside $\No$. Thus the positive square-root
and odd-degree-root tests hold in the intersection. Strict containment is
shown in the next example.
\end{proof}
If $j$ is proper, its finite iterated images are strictly decreasing:
$j^n(\No)=j^{n+1}(\No)$ would, by injectivity of $j^n$, imply
$\No=j(\No)$. Nevertheless the hierarchy stabilizes at the countable
intersection because $j(K_\infty)=K_\infty$. No stronger transfinite
recursion is needed for this stabilization.

For a pure monomial lift,
\begin{equation}\label{eq:hcore}
 K_\infty(H_\tau)=\NF{\bigcap_{n<\omega}\tau^n(\No)}.
\end{equation}
Indeed, membership is exactly the requirement that every exponent belong
to every iterated exponent image.

For $f_-$, direct iteration gives
\[
 f_-^n(a)=\frac{a}{1-na}\quad(a<0),\qquad
 f_-^n(\No)=(-1/n,\infty)\quad(n\ge1).
\]
Let
\[
 C=\bigcap_{n\ge1}(-1/n,\infty),\qquad G=\NF{C}.
\]
The class $C$ consists of nonnegative surreals together with negative
real-infinitesimals. Although $C$ itself need not be an additive group,
$G$ is an additive subgroup: supports in any fixed class are closed under
addition and negation of series. The formulas for the lifts give
\begin{equation}\label{eq:Jcore}
 K_\infty(J_-)=\NF{G}.
\end{equation}
This core is larger than $\NF{B}=\Fix(J_-)$. For example, the exponent
$a=\omega^{-1/\omega}$ lies in $G\setminus B$. Hence $\omega^a$ lies in
the core but is not fixed. A map can therefore act reversibly and
nontrivially on its core while being proper on the entire class.

\section{Topology: cofinality, discrete copies, and closed images}\label{sec:topology}

Throughout this section, the order topology is understood through interval
neighborhoods in the \emph{full} surreal class. We do not form a set of all
open subclasses. At a universe cutoff, the internally available
neighborhoods may be fewer, so an external ambient subspace topology can
be strictly finer than an internal order topology.

\subsection{Why set-indexed limits are misleading}
\begin{proposition}[Set-sized subsets are closed and discrete]
\label{prop:setdiscrete}
Every set $A\subseteq\No$ is closed and discrete in the full surreal order
topology. Every convergent set-indexed net is eventually constant.
\end{proposition}
\begin{proof}
Fix $x\in\No$. The nonzero distances $|x-a|$ for $a\in A\setminus\{x\}$
form a set of positive surreals. By interpolating the cut between $0$ and
that set, choose $\delta>0$ below every such distance. Then
$(x-\delta,x+\delta)$ meets $A$ in at most $\{x\}$. This isolates points
of $A$ and gives a neighborhood disjoint from $A$ at every point outside
$A$. For a set-indexed net, apply the same argument to its set of values
and its proposed limit. The isolating neighborhood forces eventual
constancy.
\end{proof}
The proposition is not a statement that every subclass is discrete.
The full class $\No$ is densely ordered. It says that a set-sized probe
cannot capture every infinitesimal scale of the ambient class.

\begin{corollary}[Existing set suprema are trivial]\label{cor:setsup}
If a nonempty set $A\subseteq\No$ has a supremum in $\No$, then it has a
maximum, equal to that supremum. Thus every order embedding preserves
suprema of nonempty sets whenever those suprema exist in $\No$.
\end{corollary}
\begin{proof}
If $u$ is a supremum not attained in $A$, choose $\delta>0$ below all
positive distances $u-a$. Then $u-\delta$ is still an upper bound, a
contradiction. A maximum is preserved by any order embedding.
\end{proof}
Accordingly, ``preserves all existing set suprema'' is far weaker here
than ambient simplest-cut preservation. A simplest surreal between two
sets is not generally their numerical supremum or infimum.

\subsection{The field-embedding dichotomy}
For a general order embedding, continuity at $x$ requires its images of
the left and right neighborhoods to approach $jx$ from the corresponding
sides. Cofinality of the entire image says nothing by itself about this
local condition, as $D_m$ shows. Field structure changes the answer.

\begin{theorem}[Cofinality is exactly continuity for field embeddings]
\label{thm:topdichotomy}
For an ordered-field self-embedding $j$ of $\No$, the following are
equivalent:
\begin{enumerate}[label=(\roman*)]
\item $j(\No)$ is cofinal in $\No$;
\item $j$ is continuous at $0$;
\item $j$ is continuous everywhere;
\item $j$ is a homeomorphism onto its image with the ambient subspace topology.
\end{enumerate}
If these conditions fail, the image is discrete in $\No$ and $j$ is
nowhere continuous.
\end{theorem}
\begin{proof}
A cofinal subfield has positive elements arbitrarily close to zero, by
inversion. Given $\varepsilon>0$, choose $\delta>0$ in the domain with
$0<j\delta<\varepsilon$. Then $|h|<\delta$ implies
$|jh|<\varepsilon$, proving continuity at zero. Additivity translates this
to every point.

Conversely, continuity at zero provides arbitrarily small positive image
elements: for a neighborhood of radius $\varepsilon$, use a positive
$h<\delta$ in a source neighborhood mapped into it. Inverses of these
positive image elements are cofinal.

The inverse from the image is always continuous: a source interval
$(x-\delta,x+\delta)$ is the inverse image, within the range, of the ambient
interval $(jx-j\delta,jx+j\delta)$. Thus continuity of $j$ makes it a
topological embedding.

If the subfield image $K$ is not cofinal, choose $B>0$ strictly above all
of $K$. For $0<k\in K$, its inverse also belongs to $K$, so
$k^{-1}<B$ and $k>B^{-1}$. Distinct points of $K$ are therefore separated
by more than $B^{-1}$. The image is discrete; continuity would isolate a
point of the densely ordered domain, which is impossible.
\end{proof}

\begin{theorem}[Monomial images are closed]\label{thm:closedimage}
For any ordered additive subgroup $G\subseteq\No$, the subfield
$\NF{G}$ is closed in the ambient order topology. If it is proper, it is
nowhere dense. In particular every proper image in
\cref{thm:weighted} is closed and nowhere dense.
\end{theorem}
\begin{proof}
Take $x\notin\NF{G}$. Its support has a largest exponent $a$ outside $G$.
Let $p$ be the truncation of $x$ to exponents strictly greater than $a$;
then $p\in\NF{G}$. Write $r=[x]_a\ne0$.
For any $y\in\NF{G}$, either $y-p$ has a nonzero term above $a$, in which
case that term dominates $x-y$, or the leading term of $x-y$ at exponent
$a$ is exactly $r\omega^a$. In both cases
\[
 |x-y|>\tfrac12|r|\omega^a.
\]
Hence a neighborhood of $x$ misses the subfield. This proves closedness.
A proper subfield has no interior by \cref{prop:convexfield}, so a closed
proper subfield is nowhere dense.
\end{proof}

\begin{corollary}[Strengthened stabilizer conclusion]\label{cor:strongstabilizer}
The proper embeddings in \cref{thm:stabilizer} are continuous and are
homeomorphisms onto closed nowhere-dense subfields. The embeddings of
\cref{thm:boundedcopy} have discrete images and are nowhere continuous.
Both families preserve all admissible Hahn sums.
\end{corollary}
There is therefore no implication from strong Hahn summability to
continuity in the ambient order topology. Nor does a continuous cofinal
field embedding have to have a dense image.

\subsection{An everywhere-zero derivative of a nonconstant automorphism}
There are two different uses of ``derivation'' in surreal mathematics:
an algebraic operator on field elements, and the pointwise derivative of
a function in the order topology. The latter already exhibits a striking
separation.

\begin{proposition}\label{prop:zeroderivative}
Let $c>1$ be a surreal and let $S_c=H_{a\mapsto ca}$. This is an increasing
field automorphism. As a function $\No\to\No$, its order-topological
derivative exists at every point and equals zero. Its inverse
$S_{1/c}$ has no surreal-valued derivative at any point.
\end{proposition}
\begin{proof}
Multiplication by $c$ is an ordered additive automorphism, so the field-lift
theorem applies. Additivity makes every difference quotient independent
of the base point:
\[
 \frac{S_c(x+h)-S_c(x)}h=\frac{S_c(h)}h.
\]
For $h\ne0$,
\begin{equation}\label{eq:derivval}
 \val\!\left(\frac{S_c(h)}h\right)=(c-1)\val(h).
\end{equation}
Given $\varepsilon>0$, choose $\eta$ with
$(c-1)\eta>\val(\varepsilon)$, and then take
$0<|h|<\omega^{-\eta}$. The quotient has valuation greater than
$\val(\varepsilon)$, so its absolute value is below $\varepsilon$.
This proves derivative zero by the full $\varepsilon$--$\delta$ definition.
For $S_{1/c}$ the coefficient $(1/c-1)$ is negative. As $h$ becomes small
through monomials of arbitrarily large valuation, the positive quotient
exceeds every prescribed positive surreal. It cannot have a surreal limit.
\end{proof}
The familiar real mean-value conclusion ``derivative zero implies
constant'' requires hypotheses unavailable here. A set-indexed sequence
of $h$'s is not sufficient to test the limit, by
\cref{prop:setdiscrete}. This example is not a statement about the
Berarducci--Mantova derivation, considered in \cref{sec:differential}.

\section{Valuation, leading terms, and the \texorpdfstring{$\omega$}{omega}-map}\label{sec:valuation}

\subsection{Three different valuation constraints}
For a field embedding $j$, distinguish:
\begin{align*}
\text{comparison preservation: }&
 \val x<\val y\iff\val(jx)<\val(jy);\\
\text{pointwise valuation fixation: }&
 \val(jx)=\val x;\\
\text{leading-term fixation: }&
 \operatorname{lt}(jx)=\operatorname{lt}(x).
\end{align*}
The last implies the second, which implies the first. The first is
automatic for surreal field embeddings.

\begin{proposition}\label{prop:valueaction}
Every field self-embedding of $\No$ is order preserving and induces a
unique ordered additive-group embedding $\tau_j$ of its natural value
group satisfying
\begin{equation}\label{eq:valueaction}
 \val(jx)=\tau_j(\val x)\qquad(x\ne0).
\end{equation}
The induced map on the residue field $\R$ is the identity. This does not,
by itself, say that the distinguished embedded copy of $\R$ is fixed
pointwise.
\end{proposition}
\begin{proof}
Positive elements of a real-closed field are exactly nonzero squares, so
a field embedding preserves positivity. The Archimedean comparison of two
positive elements is expressed using inequalities with ordinary integer
multiples. Since $j$ fixes the integers and reflects order, it preserves
and reflects these comparisons. Thus it induces an order embedding on
Archimedean classes. Multiplication gives additivity of the induced value
map, and surjectivity of $\val$ gives uniqueness.

The valuation ring of finite elements and its infinitesimal ideal are
preserved and reflected. The induced ordered field embedding of the
residue field $\R$ into itself fixes rational cuts, hence is the identity.
It determines the standard part of $j(r)$, not necessarily $j(r)$ itself.
\end{proof}

Within the weighted monomial class, a simple pointwise-valuation example
is obtained from $\tau=\id$ and
\[
 \chi(a)=2^{[a]_{-1}}.
\]
Coefficient extraction is additive, so this is a positive real character.
The resulting automorphism fixes $\R$ and every ordinal, preserves $\Oz$
and its floor, and fixes every valuation, but
\[
 H_{\id,\chi}\bigl(\omega^{\omega^{-1}}\bigr)
    =2\omega^{\omega^{-1}}.
\]
Pointwise value fixation is therefore not field rigidity.

\subsection{A controlled unit twist fixes even the leading term}
We now give a fully supported example outside the real-weighted monomial
class. It demonstrates why leading terms alone do not settle the field
problem.

Let $t=\omega^{-1}$, let $\beta<0$, and let
$\ell_\beta(a)=[a]_\beta\in\R$. For $r\in\R$, define the usual formal
binomial unit
\[
 (1+t)^r=\sum_{n\ge0}\binom rn t^n.
\]
These are legal Hahn series. Define
\begin{equation}\label{eq:unitlift}
 U_\beta\left(\sum_a r_a\omega^a\right)
 =\sum_a r_a\omega^a(1+t)^{\ell_\beta(a)}.
\end{equation}

\begin{proposition}[Strong leading-term stabilizers]\label{prop:unitlift}
The map $U_\beta$ is a strongly $\R$-linear field automorphism that fixes
the leading term of every nonzero surreal. It fixes $\R$ and every ordinal.
For any prescribed set $A$, $\beta<0$ can be chosen so that it also fixes
$A$ pointwise. Nevertheless it is nonidentity.
\end{proposition}
\begin{proof}
For a reverse well-ordered support $B$, the support of the proposed image
is contained in $B-\N$. This is reverse well-ordered. Each exponent has
only finitely many representations $a-n$ with $a\in B$, $n\in\N$:
otherwise $B$ would contain an infinite increasing sequence of the form
$b+n$. Thus the displayed family is summable. The same argument, together
with the local finiteness of a summable input family, proves strong
additivity.

The real binomial identity $(1+t)^{r+s}=(1+t)^r(1+t)^s$ and additivity of
$\ell_\beta$ prove multiplicativity. The term $n=0$ preserves the leading
term, while all correction terms have strictly smaller exponents.
Because $\beta<0$, we have $\ell_\beta(-1)=0$, so $U_\beta(t)=t$.
Consequently $U_\beta$ fixes every binomial unit above, by strong real
linearity. Replacing $\ell_\beta$ by $-\ell_\beta$ in the formula gives an
inverse: the two units cancel on every monomial and hence on every series.

Ordinals have no negative exponents in their own normal forms, so
$\ell_\beta$ vanishes on every ordinal exponent occurring in a Cantor
normal form. This fixes all ordinals. Taking $\beta$ below every member of
$E(A)$ similarly fixes every outer exponent of every member of $A$.
Finally,
\[
 U_\beta\bigl(\omega^{\omega^\beta}\bigr)
  =\omega^{\omega^\beta}(1+\omega^{-1})
  \ne\omega^{\omega^\beta},
\]
which proves nontriviality.
\end{proof}

This example does \emph{not} preserve omnific integrality: for $\beta<0$,
the displayed image has the negative exponent $\omega^\beta-1$, whereas
the original element belongs to $\Oz$. It therefore supplies a useful
separation from the floor-preserving monomial embeddings. It will also
fail exponential compatibility by \cref{thm:exprigidity}.

\subsection{Fixing residues need not fix real elements}
There is an explicit algebraic mechanism for moving a real by an
infinitesimal. In ZFC, choose a transcendental real $r$ and a derivation
$d:\R\to\R$ with $d(r)=1$, defined first on a transcendence basis and then
extended uniquely through the algebraic extension. Let $\Delta$ act
coefficientwise on surreal normal forms, fixing every monomial. Then
\begin{equation}\label{eq:taylorshift}
 T(x)=\sum_{n\ge0}\frac{t^n\Delta^n(x)}{n!},\qquad t=\omega^{-1},
\end{equation}
is well-defined by the same support bound $\supp(x)-\N$.
The iterated Leibniz rule proves multiplicativity. Since $\Delta(t)=0$,
the series with $-t$ is its inverse, by the binomial cancellation at each
coefficient. Thus $T$ is an ordered, strongly additive field automorphism
fixing every leading term, but $T(r)=r+t$. It is not $\R$-linear. This is
a formal operator exponential of $t\Delta$, not the Gonshor exponential
function on field elements.

\subsection{Truncation thresholds and the \texorpdfstring{$\omega$}{omega}-map}
Let $T_{>c}$ discard all normal-form terms with exponent at most $c$.
For a weighted monomial embedding the precise covariance formula is
\begin{equation}\label{eq:truncation}
 H_{\tau,\chi}T_{>c}=T_{>\tau(c)}H_{\tau,\chi}.
\end{equation}
Thus preservation of corresponding truncations does not mean commuting
with every \emph{fixed} threshold operator. Within this class, commuting
with every $T_{>c}$ forces $\tau=\id$: otherwise choose a threshold strictly
between an exponent and its image and test the monomial at that exponent.
Real weights can still remain.

The Conway map $\Omega(x)=\omega^x$ is an increasing isomorphism from the
additive group of exponents to the multiplicative monomial group. It is
not Gonshor's total exponential $\exp$. Its image is only the monomials,
whereas $\exp$ is onto all positive surreals.

\begin{theorem}[$\omega$-map and valuation rigidity]\label{thm:omegarigidity}
A field self-embedding satisfying
$j(\omega^x)=\omega^{j(x)}$ for every $x$ and
$\val(jy)=\val(y)$ for every $y\ne0$ is the identity.
\end{theorem}
\begin{proof}
For every $x$,
\[
 -jx=\val(\omega^{jx})=\val(j(\omega^x))=\val(\omega^x)=-x.
\]
No surjectivity or strong additivity is used.
\end{proof}
For a general field embedding commuting with $\Omega$, the same calculation
shows $\tau_j(x)=j(x)$. In the pure monomial class $j=H_\tau$, compatibility
is therefore exactly the functional equation
\begin{equation}\label{eq:omegafunctional}
 \tau=H_\tau.
\end{equation}
Solving this equation, and then treating nonmonomial embeddings, is a
substantial separate problem. The ordinal-fixing maps $J_-$ do not solve
it: at $x=\omega^{-1}$, their exponent action and field action differ.

\section{Exponential rigidity from finite displacement}\label{sec:exponential}

This section separates three claims: a finite identity about arbitrary
field endomorphisms, a profile theorem for ordered exponential fields,
and its application to the actual surreal exponential. The first two
have generic counterparts in the inspected repository source
\cite{RepoExp,RepoProfile}; the third still needs its formal bridge in
that development.

\subsection{The finite engine}
\begin{lemma}[Displacement amplification]\label{lem:amplify}
Let $F$ be an ordered field and $\sigma:F\to F$ a unital field
endomorphism. If $D(a)=\sigma(a)-a\ne0$, then for every $B>0$ the element
\begin{equation}\label{eq:probe}
 b=\frac{2B(1+|\sigma(a)|)}{|D(a)|}
\end{equation}
satisfies
\[
 \max\{|D(b)|,|D(ab)|\}>B.
\]
Consequently every nonidentity field endomorphism has displacement image
cofinal and coinitial in $F$.
\end{lemma}
\begin{proof}
The exact identity is the twisted Leibniz rule
\begin{equation}\label{eq:twisted}
 D(ab)=\sigma(a)D(b)+bD(a),
\end{equation}
not $aD(b)+bD(a)$. If both displayed displacements were at most $B$ in
absolute value, then
\[
 |bD(a)|=|D(ab)-\sigma(a)D(b)|\le B+|\sigma(a)|B.
\]
But the definition of $b$ makes its left side
$2B(1+|\sigma(a)|)$, a contradiction. This proves unbounded absolute
displacement. Since $D$ is additive and $D(-x)=-D(x)$, it proves both
cofinality and coinitiality.
\end{proof}

\begin{corollary}[Bounded displacement is rigid]\label{cor:boundeddisp}
If one field element $B$ bounds $|\sigma(x)-x|$ for every $x$, then
$\sigma=\id$. More generally, if all displacements lie in a proper convex
additive subgroup of $F$, then $\sigma=\id$.
\end{corollary}
\begin{proof}
For the second statement, choose $b$ outside the proper convex subgroup.
Every element of the subgroup has absolute value less than $|b|$;
otherwise convexity would place $b$ in it. Apply the first statement.
\end{proof}
The bound may be an infinite field element. Replacing it by an ordinary
integer bound would lose the intended surreal application. Notice also
that $A_{f_-}$ is a nonidentity \emph{additive} embedding with infinitesimal
displacement. Thus the field hypothesis cannot be dropped.

\subsection{Profiles in an ordered exponential field}
Let $E:(F,+,<)\to(F_{>0},\cdot,<)$ be an onto ordered exponential, and let
$w:F^\times\to\Gamma$ be an onto nontrivial convex valuation. Put
\[
 \nu(x)=w(E(x)).
\]
This is an additive, surjective, order-reversing map. Its kernel is a
proper convex additive subgroup of $F$.

\begin{theorem}[Profile rigidity]\label{thm:profile}
If a field endomorphism $\sigma$ satisfies $\nu(\sigma x)=\nu(x)$ for
every $x$, then $\sigma=\id$. In particular this holds if $\sigma$
commutes with $E$ and $w(\sigma y)=w(y)$ for every nonzero $y$.
\end{theorem}
\begin{proof}
The hypothesis places every $\sigma x-x$ in $\ker\nu$. This kernel is
proper and convex, so \cref{cor:boundeddisp} applies. For the last
assertion,
\[
 \nu(\sigma x)=w(E(\sigma x))=w(\sigma(E(x)))=w(E(x))=\nu(x).
\]
\end{proof}
This proof needs no strong summability, no normal forms, no fixed
coefficient field, and no surjectivity of $\sigma$.

\subsection{The actual surreal exponential}
Gonshor's exponential is an increasing isomorphism
$(\No,+)\to(\No_{>0},\cdot)$ extending the real exponential
\cite{Gonshor,vdDE,BM}. For its natural-valuation profile,
\begin{equation}\label{eq:expkernel}
 \ker(\val\circ\exp)=\mathcal O,
 \qquad
 \mathcal O=\{x: |x|<n\text{ for some }n\in\N\}.
\end{equation}
Indeed, a real bound on $x$ gives positive real upper and lower bounds on
$\exp x$. If $x$ is positive infinite, monotonicity puts $\exp x$ above
every real exponential of an integer, so it is infinite; for negative
infinite $x$, its exponential is infinitesimal. This proves
\eqref{eq:expkernel} without an explicit normal-form formula for $\exp$.
Every member of $\mathcal O$ has absolute value less than $\omega$.

\begin{theorem}[Exponential valuation rigidity for self-embeddings]
\label{thm:exprigidity}
Suppose $j:\No\to\No$ is a field embedding, $j\exp=\exp j$, and
$\val(jy)=\val(y)$ for every $y\ne0$. Then $j=\id$.
In particular, every exponential leading-term-preserving self-embedding
is the identity.
\end{theorem}
\begin{proof}
For every $x$,
\[
 \val(\exp(jx-x))
 =\val(j(\exp x))-\val(\exp x)=0.
\]
Thus $jx-x$ is finite by \eqref{eq:expkernel}, and all displacements have
the common bound $\omega$. Apply \cref{lem:amplify}.
\end{proof}
The conclusion applies to automorphisms as a special case and therefore
addresses the pointwise-leading-term condition in KKS Question~5.4.
The repository already presents the generic self-embedding implication;
the present contribution is an explicit integration of that argument
with the much broader embedding taxonomy. Correctness of the displayed
proof is separate from priority, and no exhaustive priority assessment
is claimed.

\begin{theorem}[Nontrivial exponential maps move values cofinally]
\label{thm:expvalue}
If $j$ is a nonidentity exponential field self-embedding and $\tau_j$ is
its value-group action, then
\[
 \{\tau_j(\gamma)-\gamma:\gamma\in\No\}
\]
is cofinal and coinitial in the value group. Consequently $j$ cannot
induce the identity modulo any proper convex subgroup of that group.
\end{theorem}
\begin{proof}
Let $\nu=\val\circ\exp$. Compatibility gives
\begin{equation}\label{eq:profileintertwine}
 \nu(jx-x)=\tau_j(\nu x)-\nu x.
\end{equation}
The image of the displacement $jx-x$ is cofinal and coinitial by
\cref{lem:amplify}. Applying the onto order-reversing map $\nu$ preserves
these two unboundedness properties, with their directions exchanged.
Since $\nu$ is onto, the right side ranges over precisely the asserted
set. A proper convex subgroup of an ordered group is bounded by any
positive element outside it, so it cannot contain this displacement set.
\end{proof}

For $J_-$ and, more generally, $J_{f_\beta}$, the exponent action
$A_f(x)-x$ has only negative exponents and is therefore infinitesimal.
Hence these maps cannot commute with the total exponential. This proves
failure of exponential compatibility without guessing a complicated
individual value of $\exp$. The unit twists and Taylor shifts of
\cref{sec:valuation} fail it for the even stronger reason that their value
action is exactly the identity.

\subsection{A uniqueness theorem beyond automorphisms}
A trivial kernel for an automorphism-group action implies faithfulness by
composing with an inverse. A self-embedding need not have a global inverse,
so the same reasoning is invalid for an embedding monoid. A direct
comparison supplies a useful replacement.

\begin{lemma}[Two-map bounded-displacement comparison]\label{lem:twomap}
Let $\sigma_1,\sigma_2:F\to F$ be field embeddings of an ordered field.
Suppose $\sigma_2(F)$ is cofinal and one $B>0$ bounds
$|\sigma_1(x)-\sigma_2(x)|$ for all $x$. Then $\sigma_1=\sigma_2$.
\end{lemma}
\begin{proof}
Put $d(x)=\sigma_1(x)-\sigma_2(x)$. The correct product identity is
\[
 d(ab)=\sigma_1(a)d(b)+\sigma_2(b)d(a).
\]
If $d(a)\ne0$, cofinality supplies $b$ such that
\[
 |\sigma_2(b)d(a)|>B(1+|\sigma_1(a)|).
\]
The product identity and the proposed bound on $d$ give the opposite
inequality. Thus there can be no such $a$.
\end{proof}

\begin{theorem}[Cofinal exponential embeddings have unique value lifts]
\label{thm:expunique}
Let $j_1,j_2:\No\to\No$ be exponential field embeddings with the same
value-group action. If at least one has cofinal image, then $j_1=j_2$.
\end{theorem}
\begin{proof}
For every $x$, equal value actions and exponential compatibility imply
\[
 \val(\exp(j_1x-j_2x))
 =\val(j_1(\exp x))-\val(j_2(\exp x))=0.
\]
Thus $j_1x-j_2x$ is finite and has uniform bound $\omega$.
Interchange the two maps if necessary and apply \cref{lem:twomap}.
\end{proof}
This extends the comparison mechanism from the presence of a surjective
map to the presence of a cofinal image. It does \emph{not} prove
uniqueness of arbitrary bounded exponential embeddings with the same
value action. That is a precise remaining question, rather than an
inference from a trivial kernel.

\subsection{What remains flexible in the exponential category}
The preceding theorems do not say that the whole exponential structure is
rigid. The known homogeneity theory for the surreal exponential field
allows nontrivial exponential automorphisms; the pointwise value-group
kernel is the restrictive condition \cite{KKS,EhrlichKaplan}.
Likewise, a nontrivial value action with cofinal displacement is only a
necessary condition for lifting to an exponential embedding. Equation
\eqref{eq:profileintertwine}, compatibility with the actual logarithm of
monomials, and compatibility with multiplication impose further conditions.
The double lift is a field construction, not a universal solution of
those conditions.

\section{Elementarity, definability, and other signatures}\label{sec:modeltheory}

\subsection{Field elementary does not mean set-theoretically elementary}
\begin{proposition}\label{prop:elementary}
Every field self-embedding of $\No$ is elementary in the language of
ordered fields. Every order self-embedding is elementary in the language
of dense linear orders without endpoints.
\end{proposition}
\begin{proof}
Both statements follow from quantifier elimination for the corresponding
theories. A field embedding preserves and reflects polynomial equations
and inequalities, and an order embedding preserves and reflects the
atomic order relations. Quantifier elimination reduces every fixed
formula to these data. This argument is a formula-by-formula scheme for
the proper-class structures and does not require a powerclass of types.
\end{proof}
In particular, the bounded copies from \cref{thm:boundedcopy} are elementary
subfields in the ordered-field language. The discrete image of $D_m$ is an
elementary suborder in the pure order language. First-order elementarity
is not ambient topological density, continuity, or simplicity preservation.

For the exponential language there is an additional standard input:
$(\No,+,\cdot,<,\exp)$ is an elementary extension of the real exponential
field, whose theory is model complete \cite{vdDE,BM,Wilkie}. Hence an
exponential field self-embedding is elementary in that language as well.
This does not make a non-exponential field embedding elementary in the
exponential expansion.

Adding simplicity, birthdays, or a predicate for $\Oz$ changes the
language again. For example the root and the two labeled immediate
successors are first-order definable in $(\No,<,\spref)$. A nonempty
prefix map fails to fix the root; $D_m$ fails to preserve immediate
successors. They are therefore not elementary embeddings of that expanded
structure. Their status as relation embeddings is weaker.

\subsection{O-minimal definability is a strong rigidity condition}
\begin{theorem}[Definable field self-embeddings are trivial]
\label{thm:definable}
Let $F$ be a real-closed field, possibly represented as a class, with an
o-minimal expansion. A field self-embedding whose graph is definable in
that expansion, with parameters allowed, is the identity.
\end{theorem}
\begin{proof}
Its fixed field is definable and infinite, since it contains $\Q$.
By o-minimality it contains a nonempty interval. A subfield containing an
interval is the whole field by \cref{prop:convexfield}.
\end{proof}
This applies to the pure ordered surreal field and to its standard
exponential expansion. In contrast, ``definable in set theory from the
normal-form construction'' is a much larger notion. The map $f_-$ itself
is given by a piecewise rational formula in the ordered field, but its
double lift $J_-$ cannot have a graph definable in that o-minimal field
language. Extracting all normal-form exponents changes the expressive
resources.

There is no conflict with the parameter-stabilizer theorem. Fixing a class
of named elements does not make a map definable from finitely many of them.
A single formula in a language with individual constants still uses only
finitely many constants.

\subsection{Algebraic variants and surcomplex extensions}
A nonzero ring homomorphism from a field to a field is injective and
unital: the image of $1$ is a nonzero idempotent, hence $1$. Dropping the
unital condition therefore only adds the zero homomorphism, which is not
an embedding. Additive-group embeddings and real-linear embeddings are
much less rigid, as $A_f$ shows. The $\omega$-map itself is a useful
additive-to-multiplicative embedding, but it is not a field endomorphism.

For the surcomplex field $\No(i)$, every real-surreal field embedding $j$
has two extensions
\[
 j_+(a+bi)=j(a)+j(b)i,\qquad
 j_-(a+bi)=j(a)-j(b)i.
\]
Both are field embeddings, and these are exactly the extensions of $j$,
because $i$ must map to one of the two roots of $X^2+1$. Their images are
$j(\No)(i)$; properness and surjectivity correspond to those of $j$.
These statements concern embeddings respecting the specified real form.
Arbitrary embeddings of the algebraically closed field $\No(i)$ need not
preserve that real form, so the two-extension classification is not a
classification of all surcomplex embeddings.

Finally, maps on pre-games or on particular cut presentations must first
respect equality of represented surreal values before they descend to
$\No$. A graph embedding of presentation trees, a map of sign sequences,
and a well-defined map of the quotient field are different constructions.
A proof assistant should retain this distinction rather than infer it
from similar notation.

\section{Differential constraints and composition}\label{sec:differential}

Fix a particular surreal derivation $\partial$, rather than merely the
assertion that some derivation exists. The Berarducci--Mantova construction
provides a strong derivation compatible with $\exp$, with constant field
$\R$; its distinguished simplest choice satisfies $\partial\omega=1$
\cite{BM}. A differential embedding must satisfy
\[
 \partial(jx)=j(\partial x).
\]
This is not the derivative of the function $x\mapsto jx$ in
\cref{prop:zeroderivative}.

\begin{proposition}[Two necessary differential conditions]
\label{prop:differential}
If $j$ commutes with a derivation whose constant field is $\R$, then $j$
fixes $\R$ pointwise. If also $\partial\omega=1$, then
$j(\omega)-\omega\in\R$.
\end{proposition}
\begin{proof}
For $r\in\R$, $\partial(jr)=j(0)=0$, so $jr\in\R$. The restriction
$j|_\R$ is an ordered field embedding of $\R$ into itself and fixes every
rational cut, hence is the identity. Furthermore
$\partial(j\omega-\omega)=j(1)-1=0$.
\end{proof}
These are necessary conditions, not a classification. In particular,
fixing $\omega$ alone does not prove differential compatibility for the
proper maps constructed above.

On a domain of transseries where substitution by $g$ is defined and the
usual chain rule is valid, the substitution map $C_g(f)=f\circ g$ satisfies
\[
 \partial C_g(f)=C_g(\partial f)\,\partial g.
\]
If it is to commute with the derivation, testing the variable itself
forces $\partial g=1$. A domain closed under this substitution and a
support theorem are still needed. One cannot extend this observation to
a global composition on all surreals merely by writing a formal symbol
$f\circ g$.

The operator exponential in \eqref{eq:taylorshift} illustrates the same
point in a controlled setting: a fixed negative support shift made every
operator series summable. For an arbitrary surreal derivation, neither
$\sum_n\partial^n(x)/n!$ nor a proposed inverse is automatically summable.
Differential, exponential, and Hahn-continuity conditions should therefore
be formalized independently and combined only after proving their
compatibility.

\section{Cutoffs, universes, and a concrete formalization plan}\label{sec:formal}

\subsection{Set-sized cutoff models require a smaller quantifier}
Write $\No_{<\kappa}$ for the surreals of birthday below $\kappa$.
An arbitrary ordinal cutoff need not be closed under field operations.
For an uncountable regular cardinal $\kappa$, the established bounded-Hahn
identification gives
\begin{equation}\label{eq:kappahahn}
 \No_{<\kappa}=\R((\omega^{\No_{<\kappa}}))_{<\kappa},
\end{equation}
where supports have cardinality less than $\kappa$; see van den Dries and
Ehrlich and its use in KKS Remark~2.3 \cite{vdDE,KKS}. This is the right
setting for literal set-sized embedding monoids in ZFC.

\begin{proposition}[Transfer of the explicit constructions]
\label{prop:cutofftransfer}
At such a cutoff, the Hahn lifts are valid whenever their exponent maps
are self-maps of $\No_{<\kappa}$. In particular the rational lower-tail and
upper-tail maps above yield the corresponding proper cofinal and bounded
field embeddings of $\No_{<\kappa}$. The parameter-stabilizer assertions
hold for prescribed sets $A$ of size less than $\kappa$, with parameters
chosen inside $\No_{<\kappa}$.
\end{proposition}
\begin{proof}
The rational formulas stay inside the cutoff field. Their inverse formulas
on their stated ranges also stay inside it. Relabeling a support does
not change its cardinality, so \eqref{eq:kappahahn} supplies both lifts.
For $|A|<\kappa$, regularity ensures $|E(A)|<\kappa$, since it is a union
of fewer than $\kappa$ sets, each of size less than $\kappa$, twice.
The birthdays in this footprint have a bound below $\kappa$, providing
an internal lower or upper parameter. The earlier proofs then apply.
\end{proof}
The hypothesis $|A|<\kappa$ is essential to this transfer. Externally
$\No_{<\kappa}$ itself is a set, and an embedding fixing every member of
that set is the identity. Thus the full-class theorem ``fix any prescribed
set'' must not be copied verbatim into a larger-universe view of a
set-sized cutoff.

Similarly, the topological statement that every set subset is closed and
discrete becomes the statement for subsets of size less than $\kappa$.
Larger nets and larger cuts may behave differently. The field-embedding
continuity dichotomy itself needs only ordered-field arguments and
continues to hold with the cutoff's own order topology.

\subsection{Coherent restrictions rather than a fictitious last stage}
A class map $j$ has a set-valued restriction to every bounded birthday
fragment. Consequently, for each ordinal $\alpha$ there is an ordinal
$F(\alpha)$ bounding the birthdays of all $j(x)$ with $\bday(x)<\alpha$.
There is a closed unbounded class of limit ordinals $\lambda$ satisfying
$F(\alpha)<\lambda$ for every $\alpha<\lambda$: repeatedly close a starting
ordinal under $F$ and take a countable supremum. At such $\lambda$,
$j[\No_{<\lambda}]\subseteq\No_{<\lambda}$.

This does not imply that those closure ordinals are regular cardinals or
that their birthday fragments are fields. Closure under a class map and
closure under arithmetic are separate obligations. Conversely, a family
$j_\alpha$ of stage maps defines a global map only if it is coherent on
overlaps. Unrelated embeddings at larger and larger cutoffs do not form
a global embedding merely because every stage has some embedding.

An especially safe universe formulation takes $U$-small supports, builds
the carrier in the next universe, and proves lift theorems polymorphically
in $U$. The final equivalence with sign-sequence arithmetic must be an
actual theorem, not a type alias silently treating a support representation
as the already constructed field.

\subsection{What was actually inspected in ProveIt}
All paths below are relative to \code{Algebra/SurrealNumbers/} at the
snapshot printed on the title page.

\begin{longtable}{P{68mm}P{82mm}}
\toprule
Inspected source & Relevant declarations and boundary \\
\midrule
\endhead
\code{Surreal/Foundations/}\newline
\code{SignSequence.lean}
& Universe-indexed signs, first-disagreement order, ordinal embedding,
\code{birthdays\_bounded}, \code{small\_bounded}, \code{not\_small},
\code{IsPrefix}, and \code{Simpler}. This file alone is not the field/Hahn
identification. \\
\code{Surreal/Algebra/}\newline
\code{ExponentialProfile.lean}
& Generic \code{OrderedExp}, displacement probes,
\code{eq\_id\_of\_commute}, value maps, and profile-comparison lemmas.
Its header explicitly leaves actual surreal exponential instantiation
pending. \\
\code{docs/surreal/}\newline
\code{exponential-automorphism-}\newline
\code{rigidity/README.md}
& Research claims, the twisted-product engine, source-version audit,
and the distinction between generic Lean results and the actual surreal
application. \\
\code{README.md}
& Project-level account of the sign-sequence field, normal forms,
omnific integers, and its formalization ledger. This is reported project
status, not an independent build performed for this article. \\
\bottomrule
\end{longtable}

The inspected portions of the two Lean files were their first 200 and
180 lines, respectively; these include definitions, theorem headers, and
some actual proofs. Assertions about later declarations were read from
the module's detailed header, not independently checked in their complete
proof bodies. The repository also points to
\code{Surreal/Algebra/SigmaDerivation.lean}; that dependency was not separately
compiled for this article.

\subsection{Proposed modules and proof obligations}
A practical development should follow the mathematical dependencies rather
than start with a single enormous theorem about ``all embeddings.''

\paragraph{Layer 1: the sign tree.}
Extend the existing sign-sequence API with concatenation and repetition,
then prove first-disagreement preservation, prefix reflection, birthday
rigidity, and initial-image surjectivity. The properness witnesses are
concrete omitted sign sequences. The all-ordinal statements must quantify
over \code{Ordinal.\{u\}} without assuming it is small at universe $u$.

\paragraph{Layer 2: a support-relabeling interface.}
Define a generic map on small-support Hahn series from an ordered embedding
of exponent types. Prove support transport, injectivity, strong sum
preservation and reflection, and image characterization. Only afterward
add an additive hypothesis on the exponent map and prove multiplication.
This makes it impossible to confuse $A_f$ with a field homomorphism.

\paragraph{Layer 3: the double lift and stabilizers.}
Use the proved equivalence between the actual surreal field and small
normal forms. Construct $J_f$, establish functoriality and faithfulness,
and implement the two-level footprint. Its smallness theorem is as
important as its arithmetic theorem. Prove omitted-monomial witnesses,
real and ordinal fixation, and the bounded-image inequality separately.

\paragraph{Layer 4: topology and omnific data.}
State continuity through the ordinary order topology of the universe-indexed
carrier. Keep internal small sets separate from all sets in the larger
universe. The cofinality/continuity theorem is generic to ordered fields;
closedness of monomial images uses the support API. Omnific-floor
compatibility should reuse the actual floor uniqueness theorem, including
the negative-infinitesimal correction.

\paragraph{Layer 5: exponential bridges.}
Instantiate the generic profile lemmas with the actual surreal exponential
and natural valuation. Required bridges are: positivity and multiplicative
additivity of $\exp$, monotonicity, surjectivity onto positives, the
valuation conventions, and the finite-kernel statement
\eqref{eq:expkernel}. Then add the two-map cofinal comparison lemma, which
needs only a finite product identity and an order-cofinality hypothesis.

A schematic, \emph{not compiled}, API boundary is:
\begin{verbatim}
-- Exponent order only: an additive lift.
additiveLift (f : Exponent -> Exponent) (hf : StrictMono f)

-- Exponent addition as well: a field lift.
fieldLift (tau : AdditiveOrderEmbedding Exponent Exponent)

-- Surreal-specific bridge through the normal-form equivalence.
doubleLift (f : No -> No) (hf : StrictMono f)

-- Quantify only over a small parameter family at universe u.
fixesSmallFamily (A : I -> No) [Small.{u} I]
\end{verbatim}
The type names here describe a proposed interface; they are not claims
that declarations with those exact names already exist in Mathlib or the
repository. A verified implementation should expose its universe
parameters, exact imports, and axiom audit along with its theorem names.

\section{A research agenda with precise targets}\label{sec:agenda}

The following questions are divided between classification problems,
sharpenings of the present results, and formalization tasks. Except where
an examined published version is named, they are proposed research targets,
not assertions that the entire literature has left them open. A literature
and proof audit should precede any claim of priority.

\subsection{Exponential and \texorpdfstring{$\omega$}{omega}-map embeddings}
\begin{question}[Characterize exponential value actions]
Which ordered additive embeddings $\tau:\No\to\No$ occur as the natural
value action of an exponential field self-embedding? Start with
\eqref{eq:profileintertwine}; add the logarithms of monomials and the
multiplicative identities of the prospective field map. Cofinal
value-displacement is necessary, but a sufficient compatibility criterion
has not been supplied here.
\end{question}

\begin{question}[Remove or justify the cofinality hypothesis]
Must two exponential self-embeddings with the same value action agree
when both images are bounded? The proof of \cref{thm:expunique} fails at
one precise step: a bounded image cannot amplify a fixed nonzero difference
arbitrarily. Resolve whether a different exponential argument repairs that
step or whether bounded counterexamples exist.
\end{question}

\begin{question}[Proper exponential copies]
Classify the possible ranges of proper exponential self-embeddings,
including whether a bounded range can occur for the actual Gonshor
exponential. The bounded field copies in this article are not examples
of exponential copies. A candidate must be checked against the full
exponential, not merely its infinitesimal power series.
\end{question}

\begin{question}[$\omega$-map compatibility]
Solve $\tau=H_\tau$ in the strong pure-monomial class, and then study
$\Omega$-compatible maps outside that class. Relate this to Questions~5.6
and~5.7 of the examined KKS version, concerning nontrivial
$\Omega$-automorphisms and extension of order automorphisms of generalized
$\varepsilon$-numbers. Pointwise value fixation is already excluded by
\cref{thm:omegarigidity}; it is not the remaining case.
\end{question}

\subsection{Combining tree structure with arithmetic}
\begin{question}[Field and simplicity together]
Classify field self-embeddings that preserve the sign-prefix relation.
Neither the general Hahn construction nor the prefix construction solves
this combined problem. An initial image is impossible for a proper map,
but prefix preservation need not imply initiality of the image.
\end{question}

\begin{question}[Elementary embeddings of the sign hierarchy]
Determine the proper elementary self-embeddings, if any, of
$(\No,<,\spref)$ and of an explicitly two-sorted expansion by birthdays
and ordinals. Such maps must fix the root and preserve labeled immediate
successors. Identify which limit-stage properties become first-order
visible and which remain external chain-completeness requirements.
\end{question}

\begin{question}[Weaker birthday controls]
Replace the global inequality $\bday(jx)\le\bday(x)$ by bounds on selected
stages, a cofinal family, or a prescribed ordinal growth function. Determine
when those weaker bounds force identity, surjectivity, or a specific
substructure. The ordinal addition in prefix maps and ordinal multiplication
in repetition maps provide useful sharp test cases.
\end{question}

\subsection{Hahn images, fixed fields, and dynamics}
\begin{question}[Beyond monomial images]
Which strongly additive field embeddings admit a factorization into a
pure exponent lift, a real character, and a controlled unit twist? The
unit construction \eqref{eq:unitlift} proves that support control is
essential. A general homomorphism into the unit group need not support a
well-defined strongly additive substitution.
\end{question}

\begin{question}[Prescribe the fixed field]
Which real-closed subfields of $\No$ arise as fixed fields of proper
embeddings or of automorphisms? The formula
$\NF{\Fix\tau\cap\ker\chi}$ solves a concrete subclass, and the pair
$J_-$, $J_g$ shows that identical fixed fields can occur on opposite sides
of surjectivity. Seek invariants that distinguish the maps when the fixed
field does not.
\end{question}

\begin{question}[Prescribe the reversible core]
Which pairs of real-closed fields $L\subseteq K\subseteq\No$ can occur as
$L=\Fix(j)$ and $K=K_\infty(j)$? For a single injective map the image
hierarchy already stabilizes at its countable intersection, by
\cref{prop:core}; there is no further transfinite descent to discover.
The remaining problem is to prescribe that stable core and its nontrivial
automorphism, and to identify the restrictions imposed by a Hahn lift.
\end{question}

\begin{question}[The stabilizer structure]
Study the composition and centralizers of the explicit families fixing
$A\cup\R\cup\On$. The double lift is faithful, so order-embedding
semigroup phenomena can be transported into the field category. Determine
which of these phenomena survive stronger constraints such as fixed
valuation, omnific floors, or a chosen derivation.
\end{question}

\begin{question}[Floor-preserving maps outside the monomial class]
Characterize field embeddings preserving $\Oz$ and its floor without
assuming monomial images. The unit twists above fail omnific preservation,
whereas every real-weighted monomial lift succeeds. Identify the support
conditions responsible for this boundary, and separately investigate
pointwise fixation of $\Oz$ rather than just predicate preservation.
\end{question}

\subsection{Topology, differentiation, and foundations}
\begin{question}[Closed images beyond strong monomial maps]
Which field self-embeddings have closed images in the full order topology?
All real-weighted monomial images do. Determine what happens for
non-strong embeddings and for immediate embeddings with unchanged value
group and residue field. Do not infer density merely from cofinality.
\end{question}

\begin{question}[Differential centralizers]
Classify embeddings commuting with the chosen Berarducci--Mantova
derivation, with or without exponential compatibility. The necessary
conditions of \cref{prop:differential} supply an inexpensive first filter.
Any proposed construction by composition should specify its domain,
chain rule, and summability theorem.
\end{question}

\begin{question}[A scale-sensitive calculus]
Compare full order-topological differentiation, valuation-based
asymptotics, and derivatives of formal transseries substitutions. Explain
systematically which additional regularity conditions exclude the
zero-derivative field automorphisms of \cref{prop:zeroderivative}. Ordinary
set sequences alone cannot supply the needed regularity test.
\end{question}

\begin{question}[Universe and cutoff absoluteness]
For transitive universes or cutoff fields, determine which embedding
properties are absolute: strong summability, cofinality, initiality, and
being elementary in a fixed language. New supports and new cuts in a
larger universe can change these answers. Coherent stage maps, rather
than arbitrary stagewise existence statements, should be the basic data.
\end{question}

\begin{question}[Minimal foundational strength]
Separate the axioms required for explicit Hahn lifts, global
back-and-forth, general class-valued recursions, and a literal carrier of all
class embeddings. Give translations between formula schemes, GBC with
specified recursion principles, and universe-indexed type theory. No one
of these should be advertised as automatically equivalent to the others.
\end{question}

\begin{question}[Near-term formalization]
First formalize the generic two-map comparison lemma and the
cofinality/continuity dichotomy; both are finite algebra/order arguments.
Next formalize support relabeling and the footprint theorem. Finally
complete the actual surreal exponential/valuation bridge and derive the
self-embedding rigidity corollary from the existing generic module.
These tasks have clear interfaces and independently checkable endpoints.
\end{question}

\section{Synthesis and an assumption checklist}\label{sec:conclusion}

The most informative answer is not a single yes or no to the existence of
nontrivial self-embeddings. It is the following set of separations.

\begin{overview}
\textbf{Flexibility with substantial arithmetic.}
There are proper elementary ordered-field self-embeddings that fix every
real, every ordinal, and any additional prescribed set; preserve Hahn
sums, corresponding normal-form truncations, omnific integrality and
floors; and are continuous with cofinal closed nowhere-dense image.
There are also bounded, nowhere-continuous elementary field copies fixing
any prescribed set and all reals.

\textbf{Rigidity from construction data.}
Nonincreasing birthdays force an order embedding to be the identity.
Preserving ambient simplest cuts does the same. A simplicity-initial
image of a full order copy must be the whole class. An order-and-simplicity
automorphism is the identity.

\textbf{Rigidity from exponential scale data.}
An exponential field self-embedding fixing natural-valuation values
pointwise is the identity. Nontrivial exponential embeddings must move
values cofinally in both directions. Among cofinal exponential embeddings,
the value action determines the map uniquely.
\end{overview}

For any proposed new theorem about a surreal self-embedding, check the
following before treating it as established: the exact language; set versus
class size; support admissibility; whether ``initial'' refers to signs or
normal forms; whether valuation is fixed pointwise or only up to an induced
map; whether continuity is topological or strong-Hahn; and whether an
inverse or a uniform truth predicate has been used without being available.
These checks are not peripheral bookkeeping. Several of the sharpest
mathematical distinctions in this article occur exactly at those boundaries.

\appendix
\section{Result and verification ledger}\label{app:ledger}

\begin{longtable}{P{56mm}P{47mm}P{47mm}}
\toprule
Result family & Mathematical basis & Verification status in this delivery \\
\midrule
\endhead
Sign sequences, normal forms, real closedness & Classical surreal theory;
repository carrier and field development & Imported standard foundations;
selected source inspection, not an independent build. \\
Order/simplicity examples and rigidity & First-disagreement and canonical
cut arguments & Complete proofs in this article; finite sign tests only
as regression checks. \\
Hahn lifts and weighted classification & Support transport and finite
coefficient convolution & Complete proofs in this article; no Lean
certification claimed. \\
Arbitrary-set stabilizers and bounded copies & Two-level support footprint
plus rational maps & Complete proofs, explicit bounds and omissions;
rational regression tests included. \\
Closed images and topology dichotomy & Ordered-field arguments and first
missing exponent & Complete proofs in this article. \\
Fixed fields and reversible core & Ordered finite-root argument and
support membership & Complete proofs in this article. \\
Unit twists and Taylor shifts & Neumann support control and formal
binomial identities & Proofs given; finite binomial checks are not an
infinite-support proof. \\
Displacement/profile rigidity & Repository finite engine and generic
profile framework & Conventional proof reproduced; generic source
inspected; actual surreal bridge not machine-checked here. \\
Cofinal comparison of two exponential maps & Two-map product identity and
finite-kernel bridge & Complete proof in this article; suggested new
generic formalization task. \\
Model-theoretic and differential background & Quantifier elimination,
Wilkie, surreal exponential and derivation literature & Imported theorems
are cited; consequences are proved or explicitly qualified. \\
Research agenda & Questions and proposed interfaces & No assertion that
all listed questions are certified open in the literature. \\
\bottomrule
\end{longtable}

The accompanying \code{verify.py} uses exact rational arithmetic. It checks
finite sign-order and prefix identities, the rational branch formulas and
iterates, finite Hahn relabeling identities, and binomial coefficients.
These tests are designed to catch transcription mistakes. They do not
prove statements about arbitrary ordinal-length signs, proper classes,
real coefficients, or all well-ordered supports. The mathematical proofs
are the argument for those generalizations.

\section{Source audit and reproducibility}\label{app:sources}

The repository inspection was pinned to commit
\code{b0d4af4f30b24529fec877e2ad4a5952888146e9}. The principal recent paper was
read in arXiv version \code{2509.22374v3}, whose submission history dates
that version to 23 April 2026. Version pinning is intentional: question
numbers and foundational conventions should not be silently transferred
between versions. Rendered HTML dates need not be publication or
submission dates. The relevant PDF page containing Questions~5.4 and~5.6
was also visually inspected during research.

The report uses no bundled external font files and no external figures.
The source is a single \code{article.tex}, with its bibliography included.
A standard TeX Live installation with \code{newtx}, \code{amsmath},
\code{tcolorbox}, \code{hyperref}, and \code{cleveref} can build it with
\code{latexmk -pdf article.tex}. The included build script and README
record the same procedure. This delivery contains no uncompiled Lean
file represented as a verified result; the API sketch in
\cref{sec:formal} is explicitly a design sketch.

\begingroup
\small
\setlength{\parskip}{0pt}
\begin{thebibliography}{99}
\setlength{\itemsep}{5pt}
\addcontentsline{toc}{section}{References}

\bibitem{Conway}
J. H. Conway, \emph{On Numbers and Games}, Academic Press, 1976;
second edition, A K Peters, 2001.

\bibitem{Gonshor}
H. Gonshor, \emph{An Introduction to the Theory of Surreal Numbers},
London Mathematical Society Lecture Note Series 110,
Cambridge University Press, 1986.

\bibitem{vdDE}
L. van den Dries and P. Ehrlich, Fields of surreal numbers and
exponentiation, \emph{Fundamenta Mathematicae} \textbf{167} (2001),
173--188; erratum, \textbf{168} (2001), 295--297.
\href{https://doi.org/10.4064/fm167-2-3}{doi:10.4064/fm167-2-3}.

\bibitem{BVH}
V. Bagayoko and J. van der Hoeven, Surreal substructures,
\emph{Fundamenta Mathematicae} \textbf{266} (2024), 25--96.
\href{https://arxiv.org/abs/2305.02001}{arXiv:2305.02001}.
Author version: \url{https://www.texmacs.org/joris/sss/sss.html}.

\bibitem{KKS}
E. Kaplan, L. S. Krapp, and M. Serra,
\emph{Decomposing the Automorphism Group of the Surreal Numbers},
\href{https://arxiv.org/abs/2509.22374v3}{arXiv:2509.22374v3},
version submitted 23 April 2026. Examined version: v3;
consulted 3 October 2026.

\bibitem{KS}
S. Kuhlmann and M. Serra, The automorphism group of a valued field of
generalised formal power series, \emph{Journal of Algebra}
\textbf{605} (2022), 339--376.
\href{https://doi.org/10.1016/j.jalgebra.2022.04.023}
{doi:10.1016/j.jalgebra.2022.04.023}.
\href{https://arxiv.org/abs/2107.03362}{arXiv:2107.03362}.

\bibitem{EhrlichKaplan}
P. Ehrlich and E. Kaplan, Surreal ordered exponential fields,
\emph{Journal of Symbolic Logic} \textbf{86} (2021), 1066--1115.
\href{https://arxiv.org/abs/2002.07739}{arXiv:2002.07739}.

\bibitem{BM}
A. Berarducci and V. Mantova, Surreal numbers, derivations and transseries,
\emph{Journal of the European Mathematical Society}
\textbf{20} (2018), 339--390.
\href{https://arxiv.org/abs/1503.00315}{arXiv:1503.00315}.

\bibitem{Wilkie}
A. J. Wilkie, Model completeness results for expansions of the ordered
field of real numbers by restricted Pfaffian functions and the exponential
function, \emph{Journal of the American Mathematical Society}
\textbf{9} (1996), 1051--1094.
\href{https://www.jstor.org/stable/2152916}{Stable article record}.

\bibitem{HamkinsETR}
J. D. Hamkins, Second-order transfinite recursion is equivalent to
Kelley--Morse set theory, author research exposition, 23 July 2017.
\href{https://jdh.hamkins.org/second-order-transfinite-recursion-is-equivalent-to-kelley-morse-set-theory/}{Author exposition and links to the underlying class-recursion work}.
Cited for the distinction between elementary class recursion and stronger
second-order recursion, not as a claim that they are the same principle.

\bibitem{RepoMain}
V. Reshetnikov, \emph{ProveIt: Surreal Numbers}, repository and research
collection, pinned commit
\code{b0d4af4f30b24529fec877e2ad4a5952888146e9}.
\href{https://github.com/VladimirReshetnikov/ProveIt/tree/b0d4af4f30b24529fec877e2ad4a5952888146e9/Algebra/SurrealNumbers}{Pinned surreal project on GitHub}.
Project README inspected 3 October 2026.

\bibitem{RepoSigns}
\emph{ProveIt}, \code{Surreal/Foundations/SignSequence.lean}, in the
surreal project of \cite{RepoMain}; first 200 lines inspected.
\href{https://github.com/VladimirReshetnikov/ProveIt/blob/b0d4af4f30b24529fec877e2ad4a5952888146e9/Algebra/SurrealNumbers/Surreal/Foundations/SignSequence.lean}{Pinned sign-sequence source on GitHub}.

\bibitem{RepoExp}
\emph{ProveIt}, \emph{Cofinal Displacement and Rigidity of Exponential
Automorphisms}, research report and accompanying README,
\code{docs/surreal/exponential-automorphism-rigidity/}, in \cite{RepoMain}.
Report dated 21 September 2026; unrefereed AI-assisted draft with generic
formalization recorded at the inspected snapshot.
\href{https://github.com/VladimirReshetnikov/ProveIt/tree/b0d4af4f30b24529fec877e2ad4a5952888146e9/Algebra/SurrealNumbers/docs/surreal/exponential-automorphism-rigidity}{Pinned report directory on GitHub}.

\bibitem{RepoProfile}
\emph{ProveIt}, \code{Surreal/Algebra/ExponentialProfile.lean}, in
\cite{RepoMain}; first 180 lines inspected, including the detailed scope
and pending-instantiation header.
\href{https://github.com/VladimirReshetnikov/ProveIt/blob/b0d4af4f30b24529fec877e2ad4a5952888146e9/Algebra/SurrealNumbers/Surreal/Algebra/ExponentialProfile.lean}{Pinned exponential-profile source on GitHub}.

\end{thebibliography}
\endgroup
\end{document}
